\clearpage
\section*{Problem 5}
\addcontentsline{toc}{section}{Problem 5}
\begingroup
\hypersetup{colorlinks=true,
            linkcolor=blue!50!black,
            citecolor=blue!50!black,
            urlcolor=blue!50!black}
\newcommand{\qfiveaO}{\mathcal{O}}
\newcommand{\qfiveSp}{\mathrm{Sp}}
\newcommand{\qfiveQ}{\mathbb{Q}}
\newcommand{\qfiveZ}{\mathbb{Z}}
\newcommand{\qfiveR}{\mathbb{R}}
\newcommand{\qfivegconn}{\operatorname{gconn}}
\newcommand{\qfiveconn}{\operatorname{conn}}
\newcommand{\qfiveSub}{\operatorname{Sub}}
\newcommand{\qfiveMap}{\operatorname{Map}}
\newcommand{\qfiveEP}{\widetilde{E\mathcal P}}
\newtheorem{qfivetheorem}{Theorem}
\newtheorem{qfivelemma}{Lemma}
\newtheorem{qfiveproposition}{Proposition}
\newtheorem{qfivecorollary}{Corollary}
\theoremstyle{definition}
\newtheorem{qfivedefinition}{Definition}
\newtheorem{qfiveremark}{Remark}
\title{Solution to Problem 5: The $\qfiveaO$-Slice Filtration for $N_\infty$ Operads}
\author{}
\date{}
\maketitle


\section*{Setup: transfer systems and indexing systems}

Fix a finite group $G$. Following Balchin--Barnes--Roitzheim \cite{BBR} and
Rubin \cite{Rubin}, the homotopy theory of $N_\infty$ operads on $G$ is governed
by the combinatorics of \emph{transfer systems}.

\begin{qfivedefinition}[transfer system]
A \emph{transfer system} on $G$ is a partial order, denoted $\to$, on the set
$\qfiveSub(G)$ of subgroups of $G$ satisfying:
\begin{enumerate}\setlength\itemsep{2pt}
\item it refines subgroup inclusion: if $H\to K$ then $H\subseteq K$;
\item it is conjugation invariant: if $H\to K$ and $g\in G$ then
$gHg^{-1}\to gKg^{-1}$;
\item it is closed under restriction: if $H\to K$ and $J\subseteq K$ then
$H\cap J\to J$.
\end{enumerate}
Transfer systems on $G$ form a poset under refinement.
\end{qfivedefinition}

\begin{qfivedefinition}[admissible sets, indexing system]
Let $\qfiveaO$ be a transfer system on $G$. A finite $H$-set $T=\coprod_i H/K_i$ is
\emph{admissible} for $\qfiveaO$ if $K_i\to H$ for all $i$. The collection $\qfiveaO(H)$ of
admissible $H$-sets, as $H$ varies, is the associated \emph{indexing system}; we
also denote the corresponding $N_\infty$ operad by $\qfiveaO$.
\end{qfivedefinition}

For a finite $G$-set $T$ we have the \emph{$T$-norm}
$N^T:\qfiveSp^G\to\qfiveSp^G$, built from the Hill--Hopkins--Ravenel norms \cite{HHR16}
by $N^{G/H}(E)=N^G_H i_H^\ast(E)$ and $N^{T_0\sqcup T_1}(E)=N^{T_0}(E)\otimes
N^{T_1}(E)$. We write $i_H^\ast$ for restriction to $H$ and $\Phi^H$ for the
geometric fixed points; recall the two structural identities we shall use
repeatedly:
\begin{align}
\Phi^G\circ N^{G/H} &\simeq \Phi^H\circ i_H^\ast , \label{eq:phinorm}\\
i_K^\ast\, N^{G/H}(X) &\simeq \bigotimes_{[g]\in K\backslash G/H}
   N^{K}_{K\cap gHg^{-1}}\, c_g\, i^\ast(X)
   \;=\; N^{\,i_K^\ast(G/H)}(i_K^\ast X), \label{eq:resnorm}
\end{align}
the geometric-fixed-point formula for norms \cite[Prop.~B.96]{HHR16} and the
double-coset (norm--restriction) formula \cite[Prop.~B.93]{HHR16}. Here
$i_K^\ast(G/H)=\coprod_{[g]\in K\backslash G/H} K/(K\cap gHg^{-1})$.

\begin{qfiveremark}[on \eqref{eq:phinorm} for non-normal $H$]\label{rem:phinorm-general}
We emphasize that \eqref{eq:phinorm} holds for \emph{every} subgroup
$H\subseteq G$, not merely for normal $H$. Indeed, the Hill--Hopkins--Ravenel
diagonal formula \cite[Prop.~B.96]{HHR16} asserts that for an arbitrary
$H\subseteq G$ and any $H$-spectrum $Y$ there is a natural equivalence
\[
\Phi^G\bigl(N^G_H Y\bigr)\;\simeq\;\Phi^H(Y).
\]
(The proof in \cite{HHR16} reduces, via the geometric description of $\Phi^G$ as
$\Phi^G(-)=(\,\widetilde{E\mathcal P}\otimes(-))^G$ and the diagonal of the indexed
smash product, to the elementary computation that the $G$-fixed points of the
$G$-set $G\times_H Z$ are $Z^H$; no normality is required, and the double-coset
formula \eqref{eq:resnorm} is \emph{not} needed for this identity.) Taking
$Y=i_H^\ast X$ and using $N^{G/H}=N^G_H\circ i_H^\ast$ gives exactly
$\Phi^G\circ N^{G/H}\simeq\Phi^H\circ i_H^\ast$, which is \eqref{eq:phinorm}.
\end{qfiveremark}

The following proposition isolates, with full proof, the geometric
fixed-point computation $\Phi^G(N^TS^1)=S^{|T/G|}$ that is the engine of the
forward direction (Lemma~\ref{lem:forward}) and of the geometric upper bound
(Lemma~\ref{lem:geom}(3)). Here $|T/G|$ denotes the number of $G$-orbits of the
finite $G$-set $T$.

\begin{qfiveproposition}[geometric fixed points of a $T$-norm sphere]\label{prop:phiNT}
Let $T$ be a finite $G$-set with orbit decomposition $T\cong\coprod_{i=1}^{m}G/H_i$
(so $m=|T/G|$). Then there is a natural equivalence of (nonequivariant) spectra
\[
\Phi^G\bigl(N^T S^1\bigr)\;\simeq\;S^{\,|T/G|}\;=\;S^{m}.
\]
In particular $\qfivegconn(N^T S^1)(G)=\operatorname{conn}\Phi^G(N^TS^1)=|T/G|$.
\end{qfiveproposition}

\begin{proof}
We use three facts, each justified below: (a) $\Phi^G$ is symmetric monoidal;
(b) $\Phi^K S^1\simeq S^1$ for every subgroup $K\subseteq G$; and (c) the diagonal
formula \eqref{eq:phinorm}, valid for all subgroups (Remark
\ref{rem:phinorm-general}).

\emph{Step 1: reduce to a single orbit via monoidality (a).}
The functor $\Phi^G:\qfiveSp^G\to\qfiveSp$ is symmetric monoidal: it preserves the
unit ($\Phi^G S^0\simeq S^0$) and there are natural equivalences
$\Phi^G(X\otimes Y)\simeq\Phi^G X\otimes\Phi^G Y$
\cite[Prop.~B.96, \S2.5.2]{HHR16},\cite[\S6]{HillPrimer}. By construction the
$T$-norm is multiplicative on disjoint unions,
$N^{T_0\sqcup T_1}S^1\simeq N^{T_0}S^1\otimes N^{T_1}S^1$; applying this to the orbit
decomposition $T\cong\coprod_{i=1}^{m}G/H_i$ and then $\Phi^G$,
\[
\Phi^G\bigl(N^T S^1\bigr)
\;\simeq\;\Phi^G\Bigl(\bigotimes_{i=1}^m N^{G/H_i}S^1\Bigr)
\;\simeq\;\bigotimes_{i=1}^m\Phi^G\bigl(N^{G/H_i}S^1\bigr),
\]
where the last step is the $m$-fold iterate of the monoidality equivalence in (a).

\emph{Step 2: a single orbit contributes $S^1$.}
Fix $i$ and write $H:=H_i$. By the diagonal formula \eqref{eq:phinorm} (valid for
the possibly non-normal $H$ by Remark \ref{rem:phinorm-general}),
\[
\Phi^G\bigl(N^{G/H}S^1\bigr)\;=\;\Phi^G\bigl(N^G_H\,i_H^\ast S^1\bigr)
\;\simeq\;\Phi^H\bigl(i_H^\ast S^1\bigr).
\]
Here $i_H^\ast S^1$ is the integer-graded sphere $S^1$ regarded as an
$H$-spectrum with \emph{trivial} $H$-action (the restriction of the
$G$-equivariant suspension $S^1=\Sigma S^0$ of the unit). We now verify fact (b),
$\Phi^H S^1\simeq S^1$. The trivial-action circle $S^1$ is the representation
sphere $S^V$ for $V$ the trivial $1$-dimensional real $H$-representation. The
geometric fixed points of a representation sphere are given by
$\Phi^H S^V\simeq S^{V^H}$, where $V^H$ is the fixed subspace
\cite[\S6]{HillPrimer},\cite[Prop.~B.96]{HHR16}. Since $V$ is trivial,
$V^H=V$ has real dimension $1$, so $\Phi^H S^1\simeq S^{V^H}=S^1$. Therefore
$\Phi^G(N^{G/H_i}S^1)\simeq S^1$ for each $i$.

\emph{Step 3: assemble.}
Substituting Step 2 into the conclusion of Step 1,
\[
\Phi^G\bigl(N^T S^1\bigr)\;\simeq\;\bigotimes_{i=1}^m S^1
\;\simeq\;S^{m}\;=\;S^{\,|T/G|},
\]
using $S^1\otimes S^1\simeq S^2$ and induction on $m$ ($\bigotimes_{i=1}^m
S^1\simeq S^m$). This proves the displayed equivalence; taking connectivities,
$\qfivegconn(N^TS^1)(G)=\operatorname{conn}(S^{|T/G|})=|T/G|$.
\end{proof}


\section*{The $\qfiveaO$-slice filtration}

We use equivariant localization (nullification) to construct the relative slice
tower. The relevant closure notion is the \emph{connective} (positive) one, not
a full stable/localizing one: for a set $\mathcal G$ of objects of $\qfiveSp^G$, let
$\langle\mathcal G\rangle_{\ge}$ denote the smallest full subcategory of
$\qfiveSp^G$ that contains $\mathcal G$ and is closed under
\begin{enumerate}\setlength\itemsep{1pt}
\item[(o1)] all homotopy colimits (in particular arbitrary wedges, filtered
colimits, and cofibers of maps \emph{between its objects}, hence extensions and
positive suspensions $\Sigma^{\ell}(-)$, $\ell\ge0$);
\item[(o2)] retracts;
\item[(o3)] tensors with orbit suspension spectra $G/H_+\otimes(-)$.
\end{enumerate}
We call $\langle\mathcal G\rangle_{\ge}$ the \emph{connective subcategory
generated by} $\mathcal G$. \emph{We emphasize that $\langle\mathcal
G\rangle_{\ge}$ is \textbf{not} required to be closed under desuspension}: closing
under $\Omega=\Sigma^{-1}$ as well would, since the cells generate $\qfiveSp^G$ as a
stable category, collapse the subcategory to all of $\qfiveSp^G$ and trivialize the
filtration. Thus $\langle\mathcal G\rangle_{\ge}$ is closed under positive
suspension, cofibers of maps between its objects, extensions, retracts, colimits,
and orbit tensors, but \emph{not} under desuspension. This is the standard
generation used for slice-type ``$\ge n$'' subcategories.

\begin{qfivedefinition}[$\qfiveaO$-slice connective spectra]
For an indexing system $\qfiveaO$, let $\tau^\qfiveaO_{\ge n}$ be the equivariant
connective subcategory $\langle\,\cdot\,\rangle_{\ge}$ of $\qfiveSp^G$ generated by
\[
\bigl\{\, G_+\otimes_H N^T S^1 \;\bigm|\; H\subseteq G,\ T\in\qfiveaO(H),\ |T|\ge n \,\bigr\}.
\]
This is the category of \emph{$\qfiveaO$-slice $n$-connective} spectra. We write
$\tau^{\qfiveaO|_H}_{\ge n}\subseteq \qfiveSp^H$ for the analogous subcategory of
$\qfiveSp^H$ built from the restricted transfer system $\qfiveaO|_H$ (Definition
\ref{def:restr} below).
\end{qfivedefinition}

\begin{qfiveremark}
There is an equivariant homeomorphism $N^T S^1\cong S^{\qfiveR\cdot T}$, the
representation sphere of the permutation representation $\qfiveR\cdot T$ of $T$. So the
$n$-connective spectra are generated by representation spheres of admissible
sets of cardinality $\ge n$. This replaces the regular-representation slice cells
$G_+\wedge_H S^{m\rho_H}$ of the complete (Hill--Hopkins--Ravenel) filtration:
in the incomplete setting one is only permitted the norms indexed by sets that
$\qfiveaO$ actually admits.
\end{qfiveremark}

\begin{qfivedefinition}
For an indexing system $\qfiveaO$ and $n\ge 0$:
\begin{itemize}\setlength\itemsep{2pt}
\item the $n$th \emph{$\qfiveaO$-slice truncation} $P^n_\qfiveaO:\qfiveSp^G_{\ge 0}\to
\qfiveSp^G_{\ge 0}$ is the nullification killing $\tau^\qfiveaO_{\ge(n+1)}$;
\item the $n$th \emph{$\qfiveaO$-slice cover} $P^\qfiveaO_n$ is the fiber of
$\mathrm{Id}\Rightarrow P^{n-1}_\qfiveaO$;
\item the $n$th \emph{$\qfiveaO$-slice} $P^n_{n,\qfiveaO}(E)$ is the fiber of
$P^n_\qfiveaO(E)\to P^{n-1}_\qfiveaO(E)$.
\end{itemize}
A $G$-spectrum $E$ is an \emph{$\qfiveaO$-$n$-slice} if $E\in\tau^\qfiveaO_{\ge n}$ and
$E\to P^n_\qfiveaO E$ is an equivalence.
\end{qfivedefinition}

\begin{qfiveremark}[suspension raises slice connectivity]\label{rem:susp}
We record a structural compatibility, used only as motivation and \emph{not} as
an input to any later proof. For a finite $G$-set $T$ let $T\sqcup\ast$ denote
$T$ together with one extra orbit $G/G$ (a single $G$-fixed point). Since
$N^{G/G}S^1\simeq S^1$ and the $T$-norm is multiplicative on disjoint unions,
\[
N^{T\sqcup\ast}S^1\;\simeq\;N^{T}S^1\otimes N^{G/G}S^1\;\simeq\;N^{T}S^1\otimes S^1
\;\simeq\;\Sigma\, N^{T}S^1 .
\]
Moreover $G/G$ is admissible for \emph{every} transfer system (reflexivity gives
$G\to G$), so $T\in\qfiveaO(G)$ implies $T\sqcup\ast\in\qfiveaO(G)$ with
$|T\sqcup\ast|=|T|+1$. Hence $\Sigma$ carries each generating cell of
$\tau^\qfiveaO_{\ge k}$ to a generating cell of $\tau^\qfiveaO_{\ge(k+1)}$, and since
$\Sigma$ is exact and commutes with the closure operations (o1)--(o3), it
restricts to a functor
\[
\Sigma:\tau^\qfiveaO_{\ge k}\longrightarrow\tau^\qfiveaO_{\ge(k+1)} .
\]
This is a one-sided \emph{containment} of subcategories, $\Sigma\,\tau^\qfiveaO_{\ge
k}\subseteq\tau^\qfiveaO_{\ge(k+1)}$; we make \emph{no} claim of an equality of
subcategories, of an on-the-nose suspension identification beyond the
equivalence displayed above, nor of discreteness of the slice categories. (Slice
categories of spectra are stable $\infty$-categories with nontrivial mapping
spectra and are in general \emph{not} discrete; no such statement is used
anywhere below.) Nothing in the convergence argument or in the proof of the Main
Theorem invokes this remark.
\end{qfiveremark}

\subsection*{Convergence of the slice tower}

We now prove that the $\qfiveaO$-slice tower of a connective $G$-spectrum
converges. We isolate the two structural inputs as lemmas. The first is the
standard fact that geometric fixed points detect connectivity; the second
records that the slice cover of a connective spectrum stays in the connective
subcategory. The connectivity bound that drives everything is supplied by the
forward direction of the Main Theorem (Lemma~\ref{lem:forward}), proved below
and used here only logically downstream of its own (independent) proof.

\begin{qfivelemma}[geometric fixed points detect connectivity]\label{lem:detect-conn}
For $X\in\qfiveSp^G$ and $d\in\qfiveZ\cup\{+\infty\}$ the following are equivalent:
\begin{enumerate}\setlength\itemsep{2pt}
\item $\qfivegconn(X)(K)\ge d$ for every subgroup $K\subseteq G$, i.e.\ the
nonequivariant spectrum $\Phi^K X$ is $(d-1)$-connected for all $K$;
\item the genuine fixed-point spectrum $X^H$ is $(d-1)$-connected
(i.e.\ $\pi_j(X^H)=0$ for all $j<d$) for every $H\subseteq G$.
\end{enumerate}
More precisely, for every $H\subseteq G$,
\begin{equation}\label{eq:conn-min}
\operatorname{conn}\bigl(X^H\bigr)\;=\;\min_{K\subseteq H}\operatorname{conn}\bigl(\Phi^K X\bigr)
\;=\;\min_{K\subseteq H}\qfivegconn(X)(K),
\end{equation}
with the conventions $\min$ over a set with a $+\infty$ entry behaving in the
usual way ($+\infty$ being the connectivity of the zero spectrum).
\end{qfivelemma}

\begin{proof}
We prove \eqref{eq:conn-min} by induction on $|H|$; the equivalence of (1) and
(2) is the special case obtained by intersecting \eqref{eq:conn-min} over all
$H$ (the right-hand sides of \eqref{eq:conn-min} range over all $K\subseteq G$
as $H$ does).

For $|H|=1$ we have $X^{e}\simeq\Phi^{e}X$ (the underlying spectrum), so both
sides of \eqref{eq:conn-min} equal $\operatorname{conn}(\Phi^{e}X)$.

Assume \eqref{eq:conn-min} for all subgroups of order $<|H|$. Work in $\qfiveSp^H$
with $Y:=i_H^\ast X$. Let $\mathcal P_H$ be the family of \emph{proper}
subgroups of $H$, with universal space $E\mathcal P_{H}$ and cofiber
$\widetilde{E\mathcal P_H}=\operatorname{cofib}(E\mathcal P_{H+}\to S^0)$ in
$\qfiveSp^H$. Smashing the isotropy-separation cofiber sequence with $Y$ and
taking genuine $H$-fixed points gives a cofiber sequence of nonequivariant
spectra
\[
\bigl(E\mathcal P_{H+}\otimes Y\bigr)^H\;\longrightarrow\;Y^H\;\longrightarrow\;
\bigl(\widetilde{E\mathcal P_H}\otimes Y\bigr)^H .
\]
By the defining property of geometric fixed points,
$\bigl(\widetilde{E\mathcal P_H}\otimes Y\bigr)^H\simeq\Phi^H Y=\Phi^H X$
\cite[\S2.5.2]{HHR16},\cite[\S6]{HillPrimer}. For the left-hand term,
$E\mathcal P_H$ is an $H$-CW complex all of whose cells are of the form
$H/K\times D^m$ with $K\in\mathcal P_H$ proper; hence $E\mathcal P_{H+}\otimes Y$
is built (as a filtered colimit of finite extensions) from suspensions of
$H/K_+\otimes Y$ with $K\subsetneq H$, and
\[
\bigl(H/K_+\otimes Y\bigr)^H\;\simeq\;Y^K\;=\;(i_K^\ast X)^K\;=\;X^K
\]
by the Wirthm\"uller/induction--coinduction adjunction
$(H/K_+\otimes(-))^H\simeq(-)^K$. Consequently
$\operatorname{conn}\bigl((E\mathcal P_{H+}\otimes Y)^H\bigr)\ge
\min_{K\subsetneq H}\operatorname{conn}(X^K)$, and by the inductive hypothesis
applied to each proper $K$,
\[
\operatorname{conn}\bigl((E\mathcal P_{H+}\otimes Y)^H\bigr)\;\ge\;
\min_{K\subsetneq H}\ \min_{K'\subseteq K}\operatorname{conn}(\Phi^{K'}X)
\;=\;\min_{K'\subsetneq H}\operatorname{conn}(\Phi^{K'}X).
\]
(The colimit/extension estimate uses only that connectivity is preserved under
suspension, wedges, filtered colimits, and is non-increasing under cofibers from
objects of the listed connectivity.) Now feed these two estimates into the
cofiber sequence above. From a cofiber sequence $A\to X^H\to B$ one has
$\operatorname{conn}(X^H)\ge\min\bigl(\operatorname{conn}(A),\operatorname{conn}(B)\bigr)$,
giving the inequality ``$\ge$'' in \eqref{eq:conn-min}:
\[
\operatorname{conn}(X^H)\;\ge\;
\min\Bigl(\min_{K'\subsetneq H}\operatorname{conn}(\Phi^{K'}X),\,\operatorname{conn}(\Phi^H X)\Bigr)
=\min_{K\subseteq H}\operatorname{conn}(\Phi^K X).
\]
This establishes the inequality we shall use, namely the implication
$(1)\Rightarrow(2)$ in the form
\begin{equation}\label{eq:conn-ge}
\operatorname{conn}(X^H)\;\ge\;\min_{K\subseteq H}\qfivegconn(X)(K)\qquad(H\subseteq G).
\end{equation}

The reverse inequality ``$\le$'' in \eqref{eq:conn-min} also holds and gives the
stated equality, but is \emph{not needed in the sequel}; we indicate it for
completeness. Reading the same cofiber sequence the other way,
$\Phi^K X\simeq\bigl(\widetilde{E\mathcal P_K}\otimes i_K^\ast X\bigr)^K$ is a
retract, up to the finite isotropy-separation filtration of $X^K$, of a spectrum
assembled from $\{X^{K'}\}_{K'\subseteq K}$; concretely, the same cofiber
sequence yields
$\operatorname{conn}(\Phi^K X)\ge\min\bigl(\operatorname{conn}(X^K),\,\operatorname{conn}((E\mathcal P_{K+}\otimes i_K^\ast X)^K)\bigr)$,
and induction on $|K|$ gives
$\min_{K\subseteq H}\operatorname{conn}(\Phi^K X)\ge\operatorname{conn}(X^H)$. This is the
standard statement that the geometric fixed-point functors $\{\Phi^K\}_K$ are
jointly conservative and detect connectivity
\cite[\S2.5.2]{HHR16},\cite[\S6]{HillPrimer}.
\end{proof}

\begin{qfivelemma}[the slice cover of a connective spectrum is connective]\label{lem:cover-conn}
Let $E$ be a connective $G$-spectrum and $n\ge 0$. Then the $\qfiveaO$-slice cover
$P^\qfiveaO_{n}E=\operatorname{fib}\bigl(E\to P^{n-1}_\qfiveaO E\bigr)$ lies in
$\tau^\qfiveaO_{\ge n}$.
\end{qfivelemma}

\begin{proof}
We exhibit the slice cover as a cellular (connective) colocalization built from
the generators of $\tau^\qfiveaO_{\ge n}$, which automatically places it in
$\tau^\qfiveaO_{\ge n}$. Let
\[
\mathcal G_n:=\bigl\{\,G_+\otimes_H N^S S^1\ \bigm|\ H\subseteq G,\ S\in\qfiveaO(H),\ |S|\ge n\,\bigr\}
\]
be the generating set of $\tau^\qfiveaO_{\ge n}$. Starting from $A_0:=0\to E$,
build a sequence $A_0\to A_1\to\cdots$ of maps $A_t\to E$ in $\qfiveSp^G$ by the
small-object/cell-attachment construction: at each stage form
\[
A_{t+1}:=\operatorname{cofib}\Bigl(\bigsqcup_{(C,\alpha)}\,C\ \xrightarrow{\ \alpha\ }\ A_t\Bigr)
\quad\text{glued over }E,
\]
where the indexing $(C,\alpha)$ ranges over all maps $\alpha\colon C\to A_t$ with
$C$ a positive suspension $\Sigma^{\ell}C'$ ($\ell\ge 0$) of a generator
$C'\in\mathcal G_n$ such that the composite $C\to A_t\to E$ is null; one attaches
a cell killing the corresponding element of
$\pi_{0}\operatorname{Map}^G(C,\operatorname{fib}(A_t\to E))$. Set
$A:=\operatorname*{colim}_t A_t$, with the induced map $a\colon A\to E$.

Each $A_t$ is constructed from $\mathcal G_n$ by iterated positive suspensions,
wedges, cofibers \emph{of maps between objects already built}, and a filtered
colimit; all of these are among the closure operations (o1)--(o3) defining
$\langle\mathcal G_n\rangle_{\ge}=\tau^\qfiveaO_{\ge n}$. Hence
$A\in\tau^\qfiveaO_{\ge n}$.

\emph{The fiber $F:=\operatorname{fib}(a\colon A\to E)$ is
$\tau^\qfiveaO_{\ge n}$-null.} Indeed, the cell-attachment construction is exactly
the small-object argument for the set $\mathcal G_n$ (closed up under positive
suspensions): at the colimit, for every generator $C\in\mathcal G_n$ and every
$\ell\ge0$ the induced map
\[
\operatorname{Map}^G(\Sigma^{\ell}C,A)\;\longrightarrow\;\operatorname{Map}^G(\Sigma^{\ell}C,E)
\]
is a weak equivalence: surjectivity on $\pi_0$ holds because any map
$\Sigma^\ell C\to E$ lifts (after finitely many stages) through some $A_t$, and
injectivity/higher vanishing holds because any null-composite $\Sigma^\ell C\to A_t\to E$
is killed by an attached cell at the next stage; this is the defining outcome of
the small-object argument. Hence $\operatorname{Map}^G(\Sigma^\ell C,F)\simeq\ast$
for all $C\in\mathcal G_n$, $\ell\ge0$. Since the class of objects $C$ with
$\operatorname{Map}^G(C,F)\simeq\ast$ is closed under the operations (o1)--(o3)
(it is closed under colimits, retracts, positive suspensions, cofibers of maps
between such objects, and orbit tensors $G/K_+\otimes(-)$, by the corresponding
adjunctions and exactness of mapping spectra), and contains $\mathcal G_n$, it
contains all of $\tau^\qfiveaO_{\ge n}=\langle\mathcal G_n\rangle_{\ge}$. Thus $F$
is $\tau^\qfiveaO_{\ge n}$-null.

\emph{Identification with the slice cover.} By definition $P^{n-1}_\qfiveaO$ is
the nullification of $\qfiveSp^G_{\ge0}$ at $\tau^\qfiveaO_{\ge n}$, characterized by
its universal property: $P^{n-1}_\qfiveaO E$ is the initial $\tau^\qfiveaO_{\ge n}$-null
object under $E$, equivalently the unique factorization
$E\to L\to(\text{null objects})$ with $L$ null and
$\operatorname{fib}(E\to L)\in\tau^\qfiveaO_{\ge n}$. The map
$E\to F[1]=\operatorname{cofib}(a)$ has these properties:
$\operatorname{cofib}(a)$ is the cofiber of $A\to E$, fits in the fiber sequence
$F\to A\to E$ rotated to $A\to E\to F[1]$, so
$\operatorname{fib}(E\to F[1])=A\in\tau^\qfiveaO_{\ge n}$; and we must check
$F[1]=\operatorname{cofib}(a)$ is $\tau^\qfiveaO_{\ge n}$-null. For
$C\in\tau^\qfiveaO_{\ge n}$, the cofiber sequence $A\to E\to\operatorname{cofib}(a)$
induces a long exact sequence on $[C,-]^G$; since $a$ induces isomorphisms
$[C,A]^G\xrightarrow{\sim}[C,E]^G$ and $[\Sigma C,A]^G\xrightarrow{\sim}[\Sigma C,E]^G$
(both $C,\Sigma C\in\tau^\qfiveaO_{\ge n}$ and
$\operatorname{Map}^G(-,F)\simeq\ast$ on $\tau^\qfiveaO_{\ge n}$), exactness forces
$[C,\operatorname{cofib}(a)]^G=0$. Hence $\operatorname{cofib}(a)$ is
$\tau^\qfiveaO_{\ge n}$-null. By the universal property of nullification,
$\operatorname{cofib}(a)\simeq P^{n-1}_\qfiveaO E$, and therefore
\[
P^\qfiveaO_{n}E=\operatorname{fib}\bigl(E\to P^{n-1}_\qfiveaO E\bigr)\simeq A\in\tau^\qfiveaO_{\ge n}. \qedhere
\]
\end{proof}

\begin{qfivetheorem}[convergence of the $\qfiveaO$-slice tower]\label{thm:converge}
Let $E$ be a connective $G$-spectrum. Then the natural map
\[
E\;\longrightarrow\;\varprojlim_{n}\,P^n_\qfiveaO E
\]
is an equivalence in $\qfiveSp^G$.
\end{qfivetheorem}

\begin{proof}
For each $n\ge0$ we have the fiber sequence
$F_n\to E\to P^n_\qfiveaO E$ with $F_n:=P^\qfiveaO_{n+1}E$. Passing to the inverse
limit over $n$ (limits commute with limits, so the limit of fiber sequences is a
fiber sequence) gives a fiber sequence
\[
\varprojlim_n F_n\;\longrightarrow\;E\;\longrightarrow\;\varprojlim_n P^n_\qfiveaO E .
\]
Hence it suffices to prove $\varprojlim_n F_n\simeq 0$.

\emph{Step 1: the fibers become highly connected.}
By Lemma~\ref{lem:cover-conn}, $F_n=P^\qfiveaO_{n+1}E\in\tau^\qfiveaO_{\ge n+1}$.
By the forward direction (Lemma~\ref{lem:forward}), every
$X\in\tau^\qfiveaO_{\ge m}$ satisfies $r_K\cdot\qfivegconn(X)(K)\ge m$ for all
$K\subseteq G$, where $r_K=[K:\chi^\qfiveaO(K)]\le|K|\le|G|$. Applying this to
$X=F_n$ and $m=n+1$,
\begin{equation}\label{eq:fiber-gconn}
\qfivegconn(F_n)(K)\;\ge\;\frac{n+1}{r_K}\;\ge\;\frac{n+1}{|G|}\qquad
\text{for every }K\subseteq G .
\end{equation}
By Lemma~\ref{lem:detect-conn} (inequality~\eqref{eq:conn-ge}), this geometric
bound upgrades to a genuine one: for every $H\subseteq G$,
\begin{equation}\label{eq:fiber-conn}
\operatorname{conn}\bigl(F_n^{\,H}\bigr)\;\ge\;\min_{K\subseteq H}\qfivegconn(F_n)(K)
\;\ge\;\frac{n+1}{|G|}\;\xrightarrow[n\to\infty]{}\;+\infty .
\end{equation}

\emph{Step 2: the homotopy limit vanishes.}
We must show $W:=\varprojlim_n F_n\simeq 0$, equivalently $\pi_j(W^H)=0$ for all
$H\subseteq G$ and all $j\in\qfiveZ$ (genuine homotopy groups are jointly
conservative). For each fixed $H$, the genuine fixed-point functor
$(-)^H\colon\qfiveSp^H\to\qfiveSp$ preserves all homotopy limits, being a right
adjoint (indeed $(-)^H=\operatorname{Map}^H(S^0,-)$). Hence
$W^H\simeq\varprojlim_n F_n^{\,H}$, and there is a Milnor exact sequence
\[
0\to{\varprojlim_n}^{1}\,\pi_{j+1}\bigl(F_n^{\,H}\bigr)\to\pi_j(W^H)\to
\varprojlim_n\,\pi_{j}\bigl(F_n^{\,H}\bigr)\to 0 .
\]
By \eqref{eq:fiber-conn}, for each fixed $j$ we have
$\pi_{j}(F_n^{\,H})=0=\pi_{j+1}(F_n^{\,H})$ as soon as
$\frac{n+1}{|G|}>j+1$, i.e.\ for all sufficiently large $n$. Thus both towers
$\{\pi_{j}(F_n^{\,H})\}_n$ and $\{\pi_{j+1}(F_n^{\,H})\}_n$ are eventually zero;
an eventually-zero tower of abelian groups satisfies $\varprojlim=0$ and
$\varprojlim^1=0$ (it is Mittag--Leffler with trivial limit). Therefore
$\pi_j(W^H)=0$ for every $j$ and every $H$, so $W\simeq 0$.

Consequently $E\xrightarrow{\ \sim\ }\varprojlim_n P^n_\qfiveaO E$, as claimed.
\end{proof}

\begin{qfiveremark}
The hypothesis that $E$ be connective (bounded below) is essential and is used
twice: in Lemma~\ref{lem:cover-conn} to realize the slice cover inside the
\emph{connective} subcategory $\tau^\qfiveaO_{\ge n+1}$ (so that
Lemma~\ref{lem:forward} applies and yields \eqref{eq:fiber-gconn}), and through
\eqref{eq:fiber-conn} to make the connectivity of the fibers tend to $+\infty$.
This is the $\qfiveaO$-adapted analogue of the convergence of the
Hill--Hopkins--Ravenel slice tower on bounded-below spectra; the role played in
\cite{HHR16} by the regular-representation slice cells
$G_+\wedge_H S^{m\rho_H}$ (whose geometric connectivity grows like $m|H|$) is
played here by the admissible norm cells $G_+\otimes_H N^S S^1$, whose geometric
connectivity grows like $|S|/r_K\to\infty$ with $|S|$. What is \emph{not} used
is any false implication ``cell connectivity $\Rightarrow$ tower convergence'':
the implication runs through the membership $F_n\in\tau^\qfiveaO_{\ge n+1}$
(Lemma~\ref{lem:cover-conn}), the uniform geometric connectivity bound on the
\emph{whole} subcategory $\tau^\qfiveaO_{\ge n+1}$ (Lemma~\ref{lem:forward}), and
the detection of genuine connectivity by geometric fixed points
(Lemma~\ref{lem:detect-conn}).
\end{qfiveremark}

\section*{Geometric fixed points and the characteristic function}

Stable equivalences in $\qfiveSp^G$ are detected on geometric fixed points
$\Phi^H$ for all $H\subseteq G$, so the connectivity of geometric fixed points
is the governing invariant.

\begin{qfivedefinition}[geometric connectivity]
For $E\in\qfiveSp^G$, let $\qfivegconn(E):\qfiveSub(G)\to\qfiveZ\cup\{\pm\infty\}$ be
$\qfivegconn(E)(H):=\qfiveconn(\Phi^H E)$.
\end{qfivedefinition}

\begin{qfivedefinition}[characteristic function]
For a transfer system $\qfiveaO$, define $\chi^\qfiveaO:\qfiveSub(G)\to\qfiveSub(G)$ by
\[
\chi^\qfiveaO(H)\;:=\;\min\{K\mid K\to H\}\;=\;\bigcap_{K\to H}K .
\]
\end{qfivedefinition}

\begin{qfivelemma}[$\chi^\qfiveaO$ is well defined as a minimum]\label{lem:chi-min}
Fix $H\subseteq G$ and let $\mathcal A_H:=\{K\mid K\to H\}$ be the set of
subgroups admissible into $H$. Then $\mathcal A_H$ is closed under finite
intersections and contains $H$; consequently $\bigcap_{K\in\mathcal A_H}K$ is the
\emph{least} element of $\mathcal A_H$. In particular
$\chi^\qfiveaO(H)=\bigcap_{K\to H}K$ lies in $\mathcal A_H$, i.e.\ the minimum is
attained and $\chi^\qfiveaO(H)\to H$.
\end{qfivelemma}

\begin{proof}
First, $H\in\mathcal A_H$: a transfer system is in particular a partial order on
$\qfiveSub(G)$, hence reflexive, so $H\to H$.

\emph{Pairwise closure.} Suppose $K_1,K_2\in\mathcal A_H$, i.e.\ $K_1\to H$ and
$K_2\to H$. By axiom (1), $K_2\subseteq H$, so we may apply the restriction
axiom (3) to the relation $K_1\to H$ with $J:=K_2\subseteq H$, obtaining
$K_1\cap K_2\to K_2$. Combining this with $K_2\to H$ and using transitivity of
the partial order $\to$ yields $K_1\cap K_2\to H$, i.e.\ $K_1\cap
K_2\in\mathcal A_H$.

\emph{Finite closure.} By induction on $r\ge 1$, any intersection
$K_1\cap\dots\cap K_r$ of members of $\mathcal A_H$ lies in $\mathcal A_H$: the
case $r=1$ is trivial, and the inductive step writes
$K_1\cap\dots\cap K_{r+1}=(K_1\cap\dots\cap K_r)\cap K_{r+1}$ and applies pairwise
closure, since $K_1\cap\dots\cap K_r\in\mathcal A_H$ by the inductive hypothesis.

\emph{The minimum is attained.} As $G$ is finite, $\qfiveSub(G)$ is finite, so
$\mathcal A_H$ is a finite (nonempty) set, say $\mathcal A_H=\{K_1,\dots,K_r\}$.
Set $K_\ast:=\bigcap_{i=1}^r K_i=\bigcap_{K\to H}K$. By finite closure
$K_\ast\in\mathcal A_H$, i.e.\ $K_\ast\to H$. Moreover $K_\ast\subseteq K$ for
every $K\in\mathcal A_H$, so $K_\ast$ is the least element of $\mathcal A_H$.
Hence $\min\{K\mid K\to H\}$ exists and equals $\bigcap_{K\to H}K=K_\ast$,
justifying the definition of $\chi^\qfiveaO(H)$.
\end{proof}

By Lemma \ref{lem:chi-min}, $\chi^\qfiveaO(H)$ is well defined as a genuine
minimum; in particular $\chi^\qfiveaO(H)\to H$ (so $\chi^\qfiveaO(H)$ is itself
admissible into $H$, which is what later arguments use) and, by axiom (1),
$\chi^\qfiveaO(H)\subseteq H$.

\begin{qfivelemma}[the admissible index]\label{lem:index}
Let $H\subseteq G$ and set $r_H:=[H:\chi^\qfiveaO(H)]$.
\begin{enumerate}\setlength\itemsep{2pt}
\item $H/\chi^\qfiveaO(H)\in\qfiveaO(H)$ and $|H/\chi^\qfiveaO(H)|=r_H$.
\item For every single-orbit admissible set $H/K\in\qfiveaO(H)$ one has $[H:K]\le r_H$.
\item Consequently, for every admissible $T\in\qfiveaO(H)$, $|T|\le r_H\cdot|T/H|$.
\end{enumerate}
\end{qfivelemma}

\begin{proof}
(1) By definition $\chi^\qfiveaO(H)\to H$, so the single orbit $H/\chi^\qfiveaO(H)$ is
admissible, of cardinality $[H:\chi^\qfiveaO(H)]=r_H$.
(2) If $H/K$ is admissible then $K\to H$, so $\chi^\qfiveaO(H)\subseteq K$ by
minimality, hence $[H:K]\le[H:\chi^\qfiveaO(H)]=r_H$.
(3) Write $T=\coprod_i H/K_i$ with $K_i\to H$. Then
$|T|=\sum_i[H:K_i]\le \sum_i r_H = r_H\cdot|T/H|$ by (2).
\end{proof}

The index $r_G=[G:\chi^\qfiveaO(G)]$ is the maximal cardinality of an admissible $G$-set
with a single orbit; this subgroup, \emph{not} $|H|$ itself, is what enters the
connectivity bound.

\section*{Main Theorem}

\begin{qfivetheorem}\label{thm:main}
Let $\qfiveaO$ be a transfer system on $G$ and let $E$ be a connective $G$-spectrum.
Then $E\in\tau^\qfiveaO_{\ge n}$ if and only if for all $H\subseteq G$,
\begin{equation}\label{eq:crit}
[H:\chi^\qfiveaO(H)]\cdot\qfivegconn(E)(H)\;\ge\;n .
\end{equation}
\end{qfivetheorem}

Equivalently, writing
$n=\min_{H\subseteq G}\,[H:\chi^\qfiveaO(H)]\cdot\qfivegconn(E)(H)$, one has
$E\in\tau^\qfiveaO_{\ge n}$; this $n$ is the \emph{$\qfiveaO$-slice connectivity} of $E$.

When $\qfiveaO$ is complete, $\chi^\qfiveaO(H)=\{e\}$ for all $H$, $[H:\chi^\qfiveaO(H)]=|H|$,
and \eqref{eq:crit} recovers the Hill--Yarnall reformulation \cite{HY18} of the
Hill--Hopkins--Ravenel slice connectivity.

\section*{Geometric isotropy and geometric Mackey functors}

Throughout, $\mathcal P$ denotes the family of \emph{proper} subgroups of $G$,
$E\mathcal P_+$ the universal space (a pointed $G$-CW-complex with
$(E\mathcal P)^H\simeq\ast$ for $H\in\mathcal P$ and $(E\mathcal P)^G=\emptyset$),
and $\qfiveEP$ the cofiber of $E\mathcal P_+\to S^0$. The two facts we use about
$\qfiveEP$ are standard \cite[\S2.5, \S6]{HillPrimer},\cite[\S2.5.2]{HHR16}:

\begin{enumerate}\setlength\itemsep{3pt}
\item[(L1)] \emph{(Smashing localization.)} The functor $L:=\qfiveEP\otimes(-)$ is a
smashing Bousfield localization of $\qfiveSp^G$: $\qfiveEP\otimes\qfiveEP\simeq\qfiveEP$,
and a $G$-spectrum $Z$ is \emph{$\qfiveEP$-local} (i.e.\ $Z\xrightarrow{\sim}\qfiveEP\otimes
Z$) if and only if $E\mathcal P_+\otimes Z\simeq 0$. For $\qfiveEP$-local $Z$ and any
$X$,
\begin{equation}\label{eq:loc-maps}
[X,Z]^G\;\cong\;[\qfiveEP\otimes X,\,Z]^G ,
\end{equation}
because $[E\mathcal P_+\otimes X,Z]^G=[E\mathcal P_+\otimes X,\qfiveEP\otimes Z]^G=
[E\mathcal P_+\otimes\qfiveEP\otimes X,Z]^G=0$ (here $E\mathcal P_+\otimes\qfiveEP\simeq 0$
by definition of $\qfiveEP$), and the cofiber sequence $E\mathcal P_+\otimes X\to X\to
\qfiveEP\otimes X$ then gives \eqref{eq:loc-maps}.
\item[(L2)] \emph{($\Phi^G$ is a monoidal localization equivalence.)} $\Phi^G$ is
symmetric monoidal and preserves all colimits, and it restricts to an equivalence
of symmetric monoidal $\infty$-categories
\[
\Phi^G:\ L\,\qfiveSp^G\;\xrightarrow{\ \sim\ }\;\qfiveSp ,
\qquad \Phi^G(\qfiveEP\otimes X)\simeq\Phi^G X ,
\]
between the full subcategory of $\qfiveEP$-local objects and ordinary spectra. In
particular, for $\qfiveEP$-local $Y,Z$,
\begin{equation}\label{eq:phi-maps}
[Y,Z]^G\;\cong\;[\Phi^G Y,\Phi^G Z]_{\qfiveSp}.
\end{equation}
\end{enumerate}

\begin{qfivedefinition}[geometric Mackey functor]\label{def:geom}
A Mackey functor $M$ for $G$ is \emph{geometric} (supported at the top, in the
sense of Greenlees--May and Hill \cite[\S6]{HillPrimer}) if
\begin{enumerate}\setlength\itemsep{2pt}
\item[(M1)] $M(G/H)=0$ for every \emph{proper} subgroup $H\subsetneq G$; and
\item[(M2)] all restriction maps $\mathrm{res}^G_H:M(G/G)\to M(G/H)$ and all
transfer maps $\mathrm{tr}^G_H:M(G/H)\to M(G/G)$ to or from proper $H$ vanish
(automatic from (M1) for $H\ne G$, and there is nothing to impose at $H=G$).
\end{enumerate}
Equivalently, $M$ is the \emph{inflation} to $G$ of the abelian group
$A:=M(G/G)$ along the projection of Mackey-functor data onto its top value; we
write $M=\underline A^{\mathrm{geo}}$ and call $A$ the \emph{geometric value} of
$M$. A geometric Mackey functor is determined up to isomorphism by its geometric
value, and every abelian group $A$ occurs.
\end{qfivedefinition}

The next proposition records the homotopy-theoretic properties of geometric
Mackey functors that the converse direction actually uses; conditions (G1) and
(G2) below are the spectrum-level reformulation of Definition \ref{def:geom}.
This resolves, in particular, the requirement that the homotopy Mackey functors
arising in the geometric part of the isotropy separation be geometric and
\emph{connective}.

\begin{qfiveproposition}[Eilenberg--MacLane spectrum of a geometric Mackey functor]\label{prop:geomEM}
Let $M=\underline A^{\mathrm{geo}}$ be the geometric Mackey functor with geometric
value $A$, and let $HM\in\qfiveSp^G$ be its (genuine) Eilenberg--MacLane
spectrum, characterized by $\underline\pi_0(HM)=M$ and $\underline\pi_k(HM)=0$
for $k\ne 0$. Then:
\begin{enumerate}\setlength\itemsep{3pt}
\item[(C)] \emph{(connectivity)} $HM$ is connective: $HM\in\tau^\qfiveaO_{\ge 0}$, and
$i_H^\ast HM\in\tau^{\qfiveaO|_H}_{\ge 0}$ for every $H\subseteq G$.
\item[(G1)] $HM$ is $\qfiveEP$-local, i.e.\ $E\mathcal P_+\otimes HM\simeq 0$.
\item[(G2)] $\Phi^G HM\simeq HA$, an ordinary Eilenberg--MacLane spectrum
concentrated in degree $0$.
\end{enumerate}
\end{qfiveproposition}

\begin{proof}
\emph{(C).} By definition $HM$ has homotopy Mackey functor concentrated in
degree $0$ (it is an Eilenberg--MacLane object for the homotopy
$t$-structure of $\qfiveSp^G$), so $\underline\pi_k(HM)=0$ for $k<0$; hence $HM$ is
connective in the standard (Postnikov) $t$-structure. Connective $G$-spectra
form precisely $\tau^\qfiveaO_{\ge 0}$: the generators $G_+\otimes_H N^TS^1$ with
$|T|\ge 0$ include all orbit spheres $G/H_+=G_+\otimes_H N^{\varnothing}S^1$ (the
$|T|=0$ cells) and their suspensions, and the connective subcategory they
generate is exactly the category of connective $G$-spectra. (Indeed, the
$|T|=0$, $\ell\ge 0$ generators $G/H_+\otimes\Sigma^\ell S^0$ already generate all
connective $G$-spectra under colimits, retracts, and orbit tensors; conversely
every generator $G_+\otimes_H N^TS^1\cong G/H_+\otimes_? S^{\bR\cdot T}$ is
connective because the permutation representation $\bR\cdot T$ has nonnegative
fixed dimension at every subgroup.) Thus $HM\in\tau^\qfiveaO_{\ge 0}$. The same
argument over $H$ (using that $i_H^\ast$ is exact and preserves the homotopy
$t$-structure, so $i_H^\ast HM$ is again connective with homotopy concentrated in
degree $0$) gives $i_H^\ast HM\in\tau^{\qfiveaO|_H}_{\ge 0}$.

\emph{(G1).} The smash $E\mathcal P_+\otimes HM$ is computed by an isotropy
spectral sequence (equivalently, by filtering $E\mathcal P_+$ by its $G$-cells,
all of which are induced from \emph{proper} subgroups). Its input is the
collection of values $\pi_\ast^H(HM)=M(G/H)$ at the proper subgroups $H\in\mathcal
P$ together with the restriction maps among them; concretely,
$E\mathcal P_+\otimes HM$ lies in the connective subcategory generated by the
spectra $G/H_+\otimes HM\simeq G_+\otimes_H i_H^\ast HM$ with $H$ proper, and the
$H$-homotopy of $i_H^\ast HM$ is the restricted Mackey functor, whose value at
$H/H$ is $M(G/H)=0$ by (M1). More precisely, since $\underline\pi_0(HM)=M$ with
$M(G/H)=0$ for all proper $H$ and all restriction maps to proper $H$ vanish
(M2), the restricted spectrum $i_H^\ast HM$ has trivial $\underline\pi_0$ at the
top of $H$ and, by the same reasoning applied to all subgroups of $H$, trivial
homotopy Mackey functor; hence $i_H^\ast HM\simeq 0$ for every proper $H$. As
$E\mathcal P_+$ is built entirely from cells induced from proper subgroups,
$E\mathcal P_+\otimes HM$ is built from $G_+\otimes_H i_H^\ast HM\simeq 0$
($H\in\mathcal P$), so $E\mathcal P_+\otimes HM\simeq 0$. Therefore $HM$ is
$\qfiveEP$-local by (L1).

\emph{(G2).} By definition $\Phi^G HM\simeq(\qfiveEP\otimes HM)^G$, the genuine
$G$-fixed points of the $\qfiveEP$-localization. Its homotopy groups are computed
by the standard geometric--fixed--point computation for Eilenberg--MacLane
spectra of Mackey functors \cite[\S6]{HillPrimer}: in degree $0$ they give the
\emph{geometric fixed points} (Brauer quotient) of $M$,
\[
\pi_0\Phi^G HM= M(G/G)\Big/\!\sum_{H\subsetneq G}\mathrm{tr}^G_H\,M(G/H)
\;=\;A/0\;=\;A,
\]
the sum of transfers being zero by (M2). For the vanishing of $\pi_k$ for $k\ne 0$ we argue cleanly from (G1), avoiding
any appeal to a Tate-type spectral sequence.

\emph{Step 1: $\Phi^G HM\simeq (HM)^G$.} By definition of geometric fixed points
as $\qfiveEP$-localized genuine fixed points,
$\Phi^G HM\simeq(\qfiveEP\otimes HM)^G$
\cite[\S2.5.2]{HHR16},\cite[\S6]{HillPrimer}. By (G1), $HM$ is $\qfiveEP$-local:
indeed $E\mathcal P_+\otimes HM\simeq 0$, so smashing $HM$ with the isotropy
separation cofiber sequence $E\mathcal P_+\to S^0\to\qfiveEP$ exhibits the unit
map $HM\to\qfiveEP\otimes HM$ as the cofiber of $E\mathcal P_+\otimes HM\to HM$,
i.e.\ as an equivalence (its fiber $E\mathcal P_+\otimes HM$ is zero). Applying
the (exact) genuine fixed-point functor $(-)^G$ to this equivalence gives
\[
\Phi^G HM\;\simeq\;(\qfiveEP\otimes HM)^G\;\simeq\;(HM)^G .
\]

\emph{Step 2: homotopy of $(HM)^G$.} The homotopy groups of the genuine fixed
points are the equivariant homotopy groups, i.e.\ the value of the homotopy
Mackey functor at $G/G$:
\[
\pi_k\bigl((HM)^G\bigr)\;=\;\pi_k^G(HM)\;=\;\underline\pi_k(HM)(G/G).
\]
By the defining property of the genuine Eilenberg--MacLane spectrum $HM$ we have
$\underline\pi_0(HM)=M$ and $\underline\pi_k(HM)=0$ for $k\ne 0$; hence
\[
\pi_k\bigl((HM)^G\bigr)\;=\;
\begin{cases}
M(G/G)=A, & k=0,\\[2pt]
0, & k\ne 0.
\end{cases}
\]
In particular $\pi_0\Phi^G HM=A$, in agreement with the Brauer-quotient
computation above (the transfers from proper subgroups vanish by (M2), so the
Brauer quotient $M(G/G)/\sum_{H\subsetneq G}\mathrm{tr}^G_H M(G/H)=A/0=A$ equals
$\underline\pi_0(HM)(G/G)$, as it must).

Combining Steps 1 and 2, $\Phi^G HM$ is a (nonequivariant) spectrum with
$\pi_0=A$ and $\pi_k=0$ for $k\ne 0$, i.e.\ $\Phi^G HM\simeq HA$, an ordinary
Eilenberg--MacLane spectrum concentrated in degree $0$ with value $A=M(G/G)$.
\end{proof}

\begin{qfiveremark}\label{rem:geom-source}
Proposition \ref{prop:geomEM} supplies, for every abelian group $A$, a
\emph{connective} $G$-spectrum $HM=H\underline A^{\mathrm{geo}}$ that is
$\qfiveEP$-local with $\Phi^G HM\simeq HA$. By (L2) it is moreover the
\emph{unique} $\qfiveEP$-local $G$-spectrum with $\Phi^G HM\simeq HA$, so
$H\underline A^{\mathrm{geo}}\simeq\Phi^{-1}(HA)$; the content of Proposition
\ref{prop:geomEM} beyond mere uniqueness is the genuine \emph{connectivity}
(C), which is what Lemma \ref{lem:geom}(2) requires and which is not visible from
the abstract description $\Phi^{-1}(HA)$ alone. These three properties
(C),(G1),(G2) are the only facts about geometric Mackey functors used below; the
notion is fully defined (Definition \ref{def:geom}) and its required properties
are established (Proposition \ref{prop:geomEM}), so no circularity remains.
\end{qfiveremark}


\subsection*{A key smashing lemma}

The following lemma is the engine that converts the geometric (cardinality)
degree into the slice (orbit-count) degree. It uses only the projection formula
and the transfer-system axioms; no $\qfiveEP$ appears, so cell structure is
preserved.

\begin{qfivelemma}[degree shift]\label{lem:shift}
Fix $H\subseteq G$ and let $C_H:=N^{H/\chi^\qfiveaO(H)}S^1\in\qfiveSp^H$, an admissible
$H$-cell with $\Phi^H C_H\simeq S^1$. Then for every $m\in\qfiveZ$,
\[
C_H\otimes\bigl(-\bigr):\quad \tau^{\qfiveaO|_H}_{\ge m}\ \longrightarrow\
\tau^{\qfiveaO|_H}_{\ge m+r_H},\qquad r_H=[H:\chi^\qfiveaO(H)].
\]
In particular $C_H^{\otimes k}\otimes X\in\tau^{\qfiveaO|_H}_{\ge m+kr_H}$ whenever
$X\in\tau^{\qfiveaO|_H}_{\ge m}$, for all $k\ge 0$.
\end{qfivelemma}

\begin{proof}
By \eqref{eq:phinorm}, $\Phi^H C_H=\Phi^{\chi^\qfiveaO(H)}i_{\chi^\qfiveaO(H)}^\ast S^1
=\Phi^{\chi^\qfiveaO(H)}S^1=S^1$, since $\Phi^K S^1=S^1$ for every $K$ ($S^1$ is a
trivial-action sphere). The functor $C_H\otimes(-)$ preserves colimits and
retracts, so it suffices to send each generator of $\tau^{\qfiveaO|_H}_{\ge m}$ into
$\tau^{\qfiveaO|_H}_{\ge m+r_H}$. A generator has the form $H_+\otimes_K N^S S^1$ with
$K\subseteq H$, $S\in\qfiveaO|_H(K)$, $|S|\ge m$. By the projection (Frobenius)
formula for the adjunction $H_+\otimes_K(-)\dashv i_K^\ast$,
\[
C_H\otimes\bigl(H_+\otimes_K N^S S^1\bigr)\;\simeq\;
H_+\otimes_K\bigl(i_K^\ast C_H\otimes N^S S^1\bigr).
\]
By the restriction formula \eqref{eq:resnorm},
$i_K^\ast C_H=N^{\,i_K^\ast(H/\chi^\qfiveaO(H))}S^1$, where
$i_K^\ast(H/\chi^\qfiveaO(H))=\coprod_{[g]\in K\backslash H/\chi^\qfiveaO(H)}
K/(K\cap g\chi^\qfiveaO(H)g^{-1})$, a $K$-set of cardinality $r_H$. Each orbit is
admissible: $\chi^\qfiveaO(H)\to H$ implies $g\chi^\qfiveaO(H)g^{-1}\to H$ by axiom (2)
(as $g\in H$ here, and we work inside $H$), hence
$K\cap g\chi^\qfiveaO(H)g^{-1}\to K$ by the restriction axiom (3) applied with
$J=K\subseteq H$. Therefore $i_K^\ast C_H$ is admissible for $\qfiveaO|_K$, and so is
the disjoint union $S':=i_K^\ast(H/\chi^\qfiveaO(H))\sqcup S$, of cardinality
$r_H+|S|\ge m+r_H$. Since $N^{S'}S^1=i_K^\ast C_H\otimes N^S S^1$, we get
\[
C_H\otimes\bigl(H_+\otimes_K N^S S^1\bigr)\simeq H_+\otimes_K N^{S'}S^1
\in\tau^{\qfiveaO|_H}_{\ge m+r_H},
\]
as required. The last statement is $k$ iterations.
\end{proof}

\subsection*{Restriction of the data}

\begin{qfivedefinition}\label{def:restr}
For $H\subseteq G$, the \emph{restricted transfer system} $\qfiveaO|_H$ on
$\qfiveSub(H)$ is defined by $K\to_{\qfiveaO|_H}K'$ iff $K\to_{\qfiveaO}K'$ (for $K,K'\subseteq
H$). The three axioms are inherited.
\end{qfivedefinition}

\begin{qfivelemma}[the data restricts]\label{lem:restr}
Let $H\subseteq G$ and $K\subseteq H$. Then:
\begin{enumerate}\setlength\itemsep{2pt}
\item $\chi^{\qfiveaO|_H}(K)=\chi^\qfiveaO(K)$, hence $[K:\chi^{\qfiveaO|_H}(K)]=[K:\chi^\qfiveaO(K)]$.
\item $\qfivegconn(i_H^\ast E)(K)=\qfivegconn(E)(K)$ for all $E\in\qfiveSp^G$.
\item If $E\in\tau^\qfiveaO_{\ge n}$ then $i_H^\ast E\in\tau^{\qfiveaO|_H}_{\ge n}$.
\item If $D\in\tau^{\qfiveaO|_H}_{\ge n}$ then $G_+\otimes_H D\in\tau^\qfiveaO_{\ge n}$.
\end{enumerate}
\end{qfivelemma}

\begin{proof}
\emph{(1).} For $K\subseteq H$, axiom (1) forces every $J$ with $J\to K$ to satisfy
$J\subseteq K\subseteq H$; thus the defining set
$\{J\mid J\to K\}=\{J\subseteq K\mid J\to K\}$ consists entirely of subgroups of $H$,
and by Definition \ref{def:restr} the relation $J\to K$ holds in $\qfiveaO|_H$ iff it
holds in $\qfiveaO$. Hence the two sets $\{J\mid J\to_{\qfiveaO}K\}$ and
$\{J\mid J\to_{\qfiveaO|_H}K\}$ are equal, and so are their (intersection-)minima:
$\chi^{\qfiveaO|_H}(K)=\chi^\qfiveaO(K)$. Taking $K=H$ gives in particular
$\chi^{\qfiveaO|_H}(H)=\chi^\qfiveaO(H)$, so the two indices appearing in the criterion
agree.

\emph{(2).} For $K\subseteq H$ the geometric fixed points $\Phi^K$ factor through
restriction: $\Phi^K\simeq\Phi^K i_K^\ast$ and $i_K^\ast=i_K^\ast i_H^\ast$ for
$K\subseteq H\subseteq G$, so $\Phi^K i_H^\ast E\simeq\Phi^K E$. Taking connectivities,
$\qfivegconn(i_H^\ast E)(K)=\qfivegconn(E)(K)$.

\emph{(3).} The restriction $i_H^\ast:\qfiveSp^G\to\qfiveSp^H$ is symmetric monoidal,
preserves all colimits and retracts, and sends $G/L_+$ to a finite coproduct of
orbits $H/M_+$; consequently it carries the connective subcategory generated by a
set $\mathcal G$ of objects into the connective subcategory generated by
$i_H^\ast\mathcal G$. It therefore suffices to show that for each generator
$G_+\otimes_L N^S S^1$ of $\tau^\qfiveaO_{\ge n}$ (so $L\subseteq G$,
$S=\coprod_i L/J_i$ with $J_i\to L$, $|S|=\sum_i[L:J_i]\ge n$) the restriction
$i_H^\ast(G_+\otimes_L N^S S^1)$ lies in $\tau^{\qfiveaO|_H}_{\ge n}$.

By the Mackey double-coset formula for the induction--restriction adjunction
$G_+\otimes_L(-)\dashv i_L^\ast$,
\[
i_H^\ast\bigl(G_+\otimes_L X\bigr)\;\simeq\;
\bigvee_{[g]\in H\backslash G/L}\,
   H_+\otimes_{M_g}\, c_g\, i_{L_g}^\ast X ,
\qquad
M_g:=H\cap gLg^{-1},\ \ L_g:=g^{-1}Hg\cap L ,
\]
where $c_g:\qfiveSp^{L_g}\xrightarrow{\sim}\qfiveSp^{M_g}$ is conjugation by $g$
(note $gL_gg^{-1}=M_g$). Apply this with $X=N^S S^1$. The norm restricts and
conjugates as a norm indexed by the correspondingly transformed set: by
\eqref{eq:resnorm} (applied to each orbit summand of $S$) together with the
analogous compatibility of $N^{(-)}$ with $c_g$,
\[
c_g\, i_{L_g}^\ast\,N^S S^1\;\simeq\; N^{\,S_g} S^1,\qquad
S_g:=c_g\, i_{L_g}^\ast(S)\ \ (\text{an }M_g\text{-set}),
\]
since $c_g i_{L_g}^\ast S^1\simeq S^1$ ($S^1$ carries trivial action). Thus
\[
i_H^\ast\bigl(G_+\otimes_L N^S S^1\bigr)\;\simeq\;
\bigvee_{[g]\in H\backslash G/L}\, H_+\otimes_{M_g}\,N^{S_g}S^1 .
\]
We claim each wedge summand is a \emph{generator} of $\tau^{\qfiveaO|_H}_{\ge n}$.

\emph{Cardinality.} Restriction and conjugation of finite sets preserve
cardinality, so $|S_g|=|S|\ge n$.

\emph{Admissibility.} It suffices to treat a single orbit $S=L/J$ with $J\to L$;
the general case follows by taking disjoint unions. Restricting the $L$-set $L/J$
to $L_g\subseteq L$ gives, by the orbit double-coset decomposition,
\[
i_{L_g}^\ast(L/J)\;=\;\coprod_{[a]\in L_g\backslash L/J} L_g/(L_g\cap aJa^{-1}) .
\]
For each $a\in L$: from $J\to L$ and conjugation invariance (axiom (2)) we get
$aJa^{-1}\to aLa^{-1}=L$; then the restriction axiom (3), applied with the
subgroup $L_g\subseteq L$, yields $L_g\cap aJa^{-1}\to L_g$. Conjugating by $g$ and
using axiom (2) once more,
\[
g\bigl(L_g\cap aJa^{-1}\bigr)g^{-1}\;\to\; gL_gg^{-1}=M_g .
\]
All groups here are subgroups of $H$ (indeed of $M_g\subseteq H$), so this is a
relation of $\qfiveaO|_H$. Hence every orbit of
$S_g=c_g i_{L_g}^\ast(L/J)=\coprod_{[a]} M_g/\,g(L_g\cap aJa^{-1})g^{-1}$ is
admissible for $\qfiveaO|_H$ at $M_g$, i.e.\ $S_g\in\qfiveaO|_H(M_g)$.

Therefore $H_+\otimes_{M_g}N^{S_g}S^1$ with $M_g\subseteq H$, $S_g\in\qfiveaO|_H(M_g)$
and $|S_g|\ge n$ is one of the defining generators of $\tau^{\qfiveaO|_H}_{\ge n}$.
The wedge of generators lies in the connective subcategory $\tau^{\qfiveaO|_H}_{\ge n}$, proving (3).

\emph{(4).} The functor $G_+\otimes_H(-):\qfiveSp^H\to\qfiveSp^G$ preserves colimits,
retracts and (by associativity of induction) tensors with orbits, so as in (3) it
suffices to treat a generator $H_+\otimes_K N^S S^1$ of $\tau^{\qfiveaO|_H}_{\ge n}$
($K\subseteq H$, $S\in\qfiveaO|_H(K)$, $|S|\ge n$). Associativity of induction gives
$G_+\otimes_H(H_+\otimes_K N^S S^1)\simeq G_+\otimes_K N^S S^1$. Here $K\subseteq H
\subseteq G$ and, by part (1) (applied with the roles of $H,K$), admissibility of an
$K$-set for $\qfiveaO|_H$ coincides with admissibility for $\qfiveaO$ (both reduce to
relations $J\to K$ among subgroups of $K\subseteq H\subseteq G$, unchanged by passing
between $\qfiveaO$, $\qfiveaO|_H$, $\qfiveaO|_K$). Thus $S\in\qfiveaO(K)$, $|S|\ge n$, and
$G_+\otimes_K N^S S^1$ is a generator of $\tau^\qfiveaO_{\ge n}$, proving (4).
\end{proof}

\begin{qfiveremark}\label{rem:reduction}
Lemma \ref{lem:restr} legitimizes the phrase ``by restriction it suffices to
treat $H=G$'' used in the forward direction (Lemma \ref{lem:forward} below).
Indeed, fix $H\subseteq G$ and put $E':=i_H^\ast E\in\qfiveSp^H$, regarded over the
group $H$ with its transfer system $\qfiveaO|_H$. By part (3), $E\in\tau^\qfiveaO_{\ge n}$
implies $E'\in\tau^{\qfiveaO|_H}_{\ge n}$; by part (1),
$[H:\chi^{\qfiveaO|_H}(H)]=[H:\chi^\qfiveaO(H)]$; and by part (2),
$\qfivegconn(E')(H)=\qfivegconn(E)(H)$. Hence the inequality
$[H:\chi^\qfiveaO(H)]\cdot\qfivegconn(E)(H)\ge n$ for $E$ at the subgroup $H$ of $G$ is
\emph{identical} to the top inequality
$[H:\chi^{\qfiveaO|_H}(H)]\cdot\qfivegconn(E')(H)\ge n$ for $E'$ at the top group $H$
of $(H,\qfiveaO|_H)$. So it suffices to prove the latter, i.e.\ to treat the case of
the ambient group; this is what the proof of Lemma \ref{lem:forward} does (there
the ambient group is renamed $G$).
\end{qfiveremark}

\begin{qfivelemma}[induced spectra]\label{lem:induced}
If $i_H^\ast E\in\tau^{\qfiveaO|_H}_{\ge n}$ for all $H\in\mathcal P$, then
$E\mathcal P_+\otimes E\in\tau^\qfiveaO_{\ge n}$.
\end{qfivelemma}

\begin{proof}
$E\mathcal P_+\otimes E$ is a (filtered) colimit of finite $G$-CW-spectra built
from cells $G/H_+\otimes\Sigma^\ell E$ with $H\in\mathcal P$, $\ell\ge 0$; hence it
lies in the connective subcategory generated by the spectra
$G/H_+\otimes E\cong G_+\otimes_H i_H^\ast E$ for $H\in\mathcal P$. By hypothesis
$i_H^\ast E\in\tau^{\qfiveaO|_H}_{\ge n}$, so $G_+\otimes_H i_H^\ast E\in\tau^\qfiveaO_{\ge n}$
by Lemma \ref{lem:restr}(4). As $\tau^\qfiveaO_{\ge n}$ is closed under (o1)--(o3), the claim follows.
\end{proof}

\section*{Proof of the Main Theorem}

\subsection*{Forward direction (necessity)}

\begin{qfivelemma}\label{lem:forward}
If $E\in\tau^\qfiveaO_{\ge n}$ then $[H:\chi^\qfiveaO(H)]\cdot\qfivegconn(E)(H)\ge n$ for all
$H\subseteq G$.
\end{qfivelemma}

\begin{proof}
We must prove, for \emph{every} subgroup $H\subseteq G$, the inequality
$r_H\cdot\qfivegconn(E)(H)\ge n$, where $r_H:=[H:\chi^\qfiveaO(H)]$. We do this one
subgroup at a time. Fix $H\subseteq G$ once and for all and put
$E':=i_H^\ast E\in\qfiveSp^H$.

\emph{Step 0: per-subgroup reduction to the top group $H$.}
By Lemma~\ref{lem:restr}(3), $E\in\tau^\qfiveaO_{\ge n}$ gives
$E'\in\tau^{\qfiveaO|_H}_{\ge n}$; by Lemma~\ref{lem:restr}(1),
$[H:\chi^{\qfiveaO|_H}(H)]=[H:\chi^\qfiveaO(H)]=r_H$; and by
Lemma~\ref{lem:restr}(2), $\qfivegconn(E')(H)=\qfivegconn(E)(H)$. Hence the target
inequality $r_H\cdot\qfivegconn(E)(H)\ge n$ is \emph{identical} to
\begin{equation}\label{eq:fwd-top}
[H:\chi^{\qfiveaO|_H}(H)]\cdot\qfiveconn\bigl(\Phi^H E'\bigr)\;\ge\;n ,
\end{equation}
an inequality about the spectrum $E'\in\tau^{\qfiveaO|_H}_{\ge n}$, the transfer
system $\qfiveaO|_H$ on $H$, and the geometric fixed points $\Phi^H$ at the
\emph{top} group $H$ of the pair $(H,\qfiveaO|_H)$. Thus it suffices to prove the
following statement, applied to $(H,\qfiveaO|_H)$ and $E'$:

\begin{quote}\emph{For every finite group $G$ with transfer system $\qfiveaO$ and
every $X\in\tau^\qfiveaO_{\ge n}$, one has
$r_G\cdot\qfiveconn(\Phi^G X)\ge n$, where $r_G=[G:\chi^\qfiveaO(G)]$.}\end{quote}

(We have renamed the ambient pair $(H,\qfiveaO|_H)$ back to $(G,\qfiveaO)$ and $E'$ to
$X$; this is legitimate because $\qfiveaO|_H$ is itself a transfer system on $H$ by
Definition~\ref{def:restr}, so the statement to be proved is of exactly the same
form, now with $H$ playing the role of the ambient group. \emph{No} information
about subgroups of $H$ other than $H$ itself is used; the entire all-subgroups
content of the lemma is carried by ranging $H$ over all subgroups in Step~0.)
We prove the displayed statement for an arbitrary $(G,\qfiveaO)$.

\emph{Step 1: $\Phi^G$ annihilates cells induced from proper subgroups.}
Let $K\subsetneq G$ be a proper subgroup and $Y\in\qfiveSp^K$. We claim
$\Phi^G(G_+\otimes_K Y)\simeq 0$. Indeed, $\Phi^G(-)\simeq(\qfiveEP\otimes(-))^G$
and $\qfiveEP$ is built from cells of isotropy $G$ only, so $\Phi^G$ kills every
spectrum induced from a proper subgroup: concretely, by the projection formula
$\qfiveEP\otimes(G_+\otimes_K Y)\simeq G_+\otimes_K\bigl(i_K^\ast\qfiveEP\otimes Y\bigr)$,
and $i_K^\ast\qfiveEP\simeq 0$ for $K\in\mathcal P$ because
$(E\mathcal P)^K\simeq\ast$ forces $i_K^\ast E\mathcal P_+\xrightarrow{\sim} i_K^\ast S^0$,
whence $i_K^\ast\qfiveEP=\operatorname{cofib}(i_K^\ast E\mathcal P_+\to i_K^\ast S^0)\simeq 0$.
Therefore $\qfiveEP\otimes(G_+\otimes_K Y)\simeq 0$, and applying $(-)^G$ gives
$\Phi^G(G_+\otimes_K Y)\simeq 0$. In particular $\Phi^G$ of any such spectrum is
$\infty$-connective.

\emph{Step 2: the connectivity bound is closed under (o1)--(o3).}
Let $c:=n/r_G\in\qfiveQ$ and consider the full subcategory
\[
\mathcal C\;:=\;\bigl\{\,X\in\qfiveSp^G \ \bigm|\ \qfiveconn(\Phi^G X)\ge c\,\bigr\}
\]
(with the convention that $\infty$-connective spectra belong to $\mathcal C$).
Since $\Phi^G:\qfiveSp^G\to\qfiveSp$ is exact and preserves all homotopy colimits and
retracts, and since ordinary connectivity $\qfiveconn(-)\ge c$ is preserved by
suspension by nonnegative amounts, by cofibers/extensions, by retracts, and by
arbitrary homotopy colimits (a colimit of $\lceil c\rceil$-connective spectra is
$\lceil c\rceil$-connective, and $\qfiveconn\ge c\iff\qfiveconn\ge\lceil c\rceil$ as
$\qfiveconn$ is integer- or $\infty$-valued), the subcategory $\mathcal C$ is closed
under the operations (o1) and (o2). It is closed under
orbit tensors (o3) $G/L_+\otimes(-)$ as well: for $L=G$ this is the identity, and for
$L\subsetneq G$ we have $G/L_+\otimes X\cong G_+\otimes_L i_L^\ast X$, which lies in
$\mathcal C$ by Step~1 (it is $\infty$-connective after $\Phi^G$). Hence
$\mathcal C$ is closed under (o1)--(o3).

\emph{Step 3: every generator of $\tau^\qfiveaO_{\ge n}$ lies in $\mathcal C$.}
The generators are $G_+\otimes_L N^S S^1$ with $L\subseteq G$, $S\in\qfiveaO(L)$,
$|S|\ge n$. If $L\subsetneq G$ is proper, then $\Phi^G(G_+\otimes_L N^S S^1)\simeq 0$
by Step~1, so the generator is $\infty$-connective after $\Phi^G$ and lies in
$\mathcal C$. If $L=G$, the generator is $N^T S^1$ for an admissible $G$-set
$T:=S\in\qfiveaO(G)$ with $|T|\ge n$; by Proposition~\ref{prop:phiNT},
$\Phi^G(N^T S^1)\simeq S^{|T/G|}$, so $\qfiveconn(\Phi^G N^T S^1)=|T/G|$, and by
Lemma~\ref{lem:index}(3),
\[
r_G\cdot|T/G|\;\ge\;|T|\;\ge\;n,\qquad\text{i.e.}\qquad |T/G|\ge n/r_G=c .
\]
Hence $N^T S^1\in\mathcal C$ as well. In all cases the generators lie in
$\mathcal C$.

\emph{Step 4: conclusion of the displayed statement.}
By Step~2, $\mathcal C$ is closed under (o1)--(o3), and by Step~3 it contains all generators of
$\tau^\qfiveaO_{\ge n}$; since $\tau^\qfiveaO_{\ge n}=\langle\,\text{generators}\,\rangle_{\ge}$ is the \emph{smallest}
subcategory containing those generators and closed under (o1)--(o3), $\tau^\qfiveaO_{\ge n}\subseteq\mathcal C$.
Thus every $X\in\tau^\qfiveaO_{\ge n}$ satisfies $\qfiveconn(\Phi^G X)\ge c=n/r_G$,
i.e.\ $r_G\cdot\qfiveconn(\Phi^G X)\ge n$. This proves the displayed statement.

Applying the displayed statement to $(H,\qfiveaO|_H)$ and $X=E'$ yields exactly
\eqref{eq:fwd-top}, hence $r_H\cdot\qfivegconn(E)(H)\ge n$. As $H\subseteq G$ was
arbitrary, the lemma is proved.
\end{proof}

\subsection*{Converse (sufficiency)}

We prove the converse by downward induction on the order of the ambient group,
phrased uniformly: for a finite group $G$ with transfer system $\qfiveaO$, let
$P(G,\qfiveaO)$ be the statement of the converse, namely
\begin{quote}
\emph{every connective $E\in\qfiveSp^G$ with $[H:\chi^\qfiveaO(H)]\cdot\qfivegconn(E)(H)\ge n$
for all $H\subseteq G$ lies in $\tau^\qfiveaO_{\ge n}$.}
\end{quote}
\medskip
\noindent\emph{Base case ($|G|=1$).} The only subgroup is $H=e$, and the unique
transfer system is trivial; $\chi^\qfiveaO(e)=e$ and $r_e=[e:e]=1$. For the trivial
group $\qfiveSp^G=\qfiveSp$, geometric fixed points $\Phi^e=\mathrm{id}$, and the only
admissible $e$-sets are the finite trivial sets $T=\{\ast\}^{\sqcup m}$ with
$N^TS^1=S^m$; hence $\tau^\qfiveaO_{\ge n}$ is the connective subcategory generated by
$\{S^m:m\ge n\}$, which is exactly the ordinary $n$-connective category
$\qfiveSp_{\ge n}$. The hypothesis $[e:\chi^\qfiveaO(e)]\cdot\qfivegconn(E)(e)\ge n$ reads
$\qfiveconn(E)\ge n$, which is precisely the statement $E\in\qfiveSp_{\ge n}=
\tau^\qfiveaO_{\ge n}$. So $P(e,\qfiveaO)$ holds. (This is the $\qfiveaO$-slice filtration
reducing to the ordinary Postnikov filtration, as it must.)

\medskip
\noindent\emph{Inductive hypothesis.} Fix $G$ with $|G|>1$ and assume that
$P(K,\qfiveaO|_K)$ holds \emph{for every proper subgroup $K\subsetneq G$}, taken with
its restricted transfer system $\qfiveaO|_K$ (Definition \ref{def:restr}). Explicitly:
for every proper $K$ and every connective $D\in\qfiveSp^K$ with
$[J:\chi^{\qfiveaO|_K}(J)]\cdot\qfivegconn(D)(J)\ge n$ for all $J\subseteq K$, one has
$D\in\tau^{\qfiveaO|_K}_{\ge n}$. We deduce $P(G,\qfiveaO)$.

\begin{qfiveremark}\label{rem:ind-link}
The induction is over the order of the ambient group, and at the inductive step
the hypothesis is applied to the \emph{restricted} datum $(K,\qfiveaO|_K)$ acting on
the $K$-spectrum $D=i_K^\ast E$, never to $E$ itself over $G$. The bridge between
the conclusion ``$i_K^\ast E\in\tau^{\qfiveaO|_K}_{\ge n}$'' it produces and the
$G$-level subcategory $\tau^\qfiveaO_{\ge n}$ is supplied entirely by Lemma
\ref{lem:restr}: part (1) identifies $\chi^{\qfiveaO|_K}=\chi^\qfiveaO$ on subgroups of
$K$, part (2) identifies $\qfivegconn(i_K^\ast E)=\qfivegconn(E)$ on subgroups of $K$
(so the hypotheses transport), and part (4) converts $\tau^{\qfiveaO|_K}_{\ge n}$ back
to $\tau^\qfiveaO_{\ge n}$ via induction $G_+\otimes_K(-)$. In particular we do
\emph{not} require any identification of $i_K^\ast(\tau^\qfiveaO_{\ge n})$ with
$\tau^{\qfiveaO|_K}_{\ge n}$: Lemma \ref{lem:induced} is stated and used in the form
``$i_K^\ast E\in\tau^{\qfiveaO|_K}_{\ge n}$ for all proper $K\Rightarrow E\mathcal
P_+\otimes E\in\tau^\qfiveaO_{\ge n}$'', matching exactly what the inductive hypothesis
delivers.
\end{qfiveremark}

The central computation is the following, which replaces the unproved equivalence
in the original Lemma 2.

\begin{qfivelemma}[geometric slices]\label{lem:geom}
Let $M$ be a geometric Mackey functor (Definition \ref{def:geom}), so that its
Eilenberg--MacLane spectrum $HM$ is connective, $\qfiveEP$-local, and has
$\Phi^G HM\simeq HA$ in degree $0$ (Proposition \ref{prop:geomEM}: properties
(C),(G1),(G2)), and let $k\ge 0$.
Set $r_G=[G:\chi^\qfiveaO(G)]$ and $C:=C_G=N^{G/\chi^\qfiveaO(G)}S^1$. Then:
\begin{enumerate}\setlength\itemsep{3pt}
\item \emph{(identity)} $\qfiveEP\otimes\Sigma^k HM\ \simeq\ \qfiveEP\otimes C^{\otimes k}\otimes HM$;
\item \emph{(membership / lower bound)} $\qfiveEP\otimes\Sigma^k HM\in\tau^\qfiveaO_{\ge k r_G}$;
\item \emph{(vanishing / upper bound)} for every admissible $T\in\qfiveaO(G)$ with
$|T|>k r_G$,
\[
[N^T S^1,\Sigma^k HM]^G\;=\;0 .
\]
\end{enumerate}
\end{qfivelemma}

\begin{proof}
\textbf{(1)} Both sides are $\qfiveEP$-local: the right side is $\qfiveEP\otimes(-)$ of
something, and the left side is $\qfiveEP\otimes\Sigma^k HM$. By (L2) it suffices to
compare their geometric fixed points. Using that $\Phi^G$ is monoidal,
\eqref{eq:phinorm}, $\Phi^G S^1=S^1$, and (G2),
\[
\Phi^G\bigl(\qfiveEP\otimes C^{\otimes k}\otimes HM\bigr)
=(\Phi^G C)^{\otimes k}\otimes\Phi^G HM
=(S^1)^{\otimes k}\otimes HA=\Sigma^k HA ,
\]
\[
\Phi^G\bigl(\qfiveEP\otimes\Sigma^k HM\bigr)=\Sigma^k\Phi^G HM=\Sigma^k HA .
\]
The two are equivalent, so by (L2) the original $\qfiveEP$-local spectra are
equivalent.

\textbf{(2)} Put $Z:=C^{\otimes k}\otimes HM\in\qfiveSp^G$. By Proposition
\ref{prop:geomEM}(C), $HM$ is connective, so $HM\in\tau^\qfiveaO_{\ge 0}$ (connective
$G$-spectra are exactly $\tau^\qfiveaO_{\ge 0}$, generated by the $|T|=0$ cells
$G/H_+\otimes\Sigma^\ell S^0$, $\ell\ge 0$). By Lemma \ref{lem:shift} (with $H=G$,
$m=0$), $Z=C^{\otimes k}\otimes HM\in\tau^\qfiveaO_{\ge k r_G}$.

Consider the isotropy-separation cofiber sequence
\[
E\mathcal P_+\otimes Z\;\longrightarrow\; Z\;\longrightarrow\;\qfiveEP\otimes Z .
\]
By part (1), $\qfiveEP\otimes Z=\qfiveEP\otimes\Sigma^k HM$, which is therefore the cofiber
of $E\mathcal P_+\otimes Z\to Z$. We show both source and target of this map lie
in $\tau^\qfiveaO_{\ge k r_G}$; as that subcategory is closed under cofibers, the
conclusion (2) follows.

We already have $Z\in\tau^\qfiveaO_{\ge k r_G}$. For $E\mathcal P_+\otimes Z$ we apply
Lemma \ref{lem:induced}: we must check $i_H^\ast Z\in\tau^{\qfiveaO|_H}_{\ge k r_G}$ for
all proper $H$. By \eqref{eq:resnorm}, $i_H^\ast C=N^{\,i_H^\ast(G/\chi^\qfiveaO(G))}S^1$
is an admissible $H$-cell of cardinality $r_G$, and the same double-coset
admissibility argument as in Lemma \ref{lem:shift} shows each of its orbits is
admissible for $\qfiveaO|_H$; moreover $\Phi^H i_H^\ast C$ need not be $S^1$, but we do
not need that. Instead observe directly that $i_H^\ast C$ is itself an admissible
$H$-cell of cardinality $r_G$, so smashing with it raises the $\qfiveaO|_H$-degree by
$r_G$: indeed $i_H^\ast C\otimes(H_+\otimes_K N^S S^1)\simeq
H_+\otimes_K(i_K^\ast i_H^\ast C\otimes N^S S^1)=H_+\otimes_K N^{S''}S^1$ with $S''$
admissible of cardinality $r_G+|S|$ (same projection-formula computation as in
Lemma \ref{lem:shift}, using $i_K^\ast i_H^\ast=i_K^\ast$ and that
$i_K^\ast i_H^\ast C$ is an admissible $K$-cell of cardinality $r_G$). Hence
$i_H^\ast C\otimes(-)$ maps $\tau^{\qfiveaO|_H}_{\ge m}$ to $\tau^{\qfiveaO|_H}_{\ge m+r_G}$.
Iterating $k$ times and using $i_H^\ast HM\in\tau^{\qfiveaO|_H}_{\ge 0}$
(Proposition \ref{prop:geomEM}(C)), we obtain
\[
i_H^\ast Z=(i_H^\ast C)^{\otimes k}\otimes i_H^\ast HM\in\tau^{\qfiveaO|_H}_{\ge k r_G}.
\]
By Lemma \ref{lem:induced}, $E\mathcal P_+\otimes Z\in\tau^\qfiveaO_{\ge k r_G}$. This
completes (2). (Note: this argument is self-contained and does not invoke the
inductive hypothesis $P(K,\qfiveaO|_K)$.)

\textbf{(3)} Let $T\in\qfiveaO(G)$ with $|T|>k r_G$. By Lemma \ref{lem:index}(3),
$|T|\le r_G\,|T/G|$, so $|T/G|>k$.

We stress that the following computation does \emph{not} assert that
$[X,\Sigma^kHM]^G$ is determined by $\Phi^GX$ alone for an arbitrary source $X$
(which is false in general, as it would require knowing all $\Phi^HX$). What we
use instead is that the \emph{target} $\Sigma^k HM$ is $\qfiveEP$-local: this lets
us first replace $X=N^TS^1$ by its localization $\qfiveEP\otimes X$ at no cost
(\eqref{eq:loc-maps}), and then invoke that $\Phi^G$ is an \emph{equivalence}
on $\qfiveEP$-local objects (L2). The passage through $\Phi^G$ is thus a
consequence of locality of the target, not an unjustified shortcut.

Since $HM$ is $\qfiveEP$-local (G1), so is
$\Sigma^k HM$; by \eqref{eq:loc-maps} and then \eqref{eq:phi-maps},
\[
[N^T S^1,\Sigma^k HM]^G
\overset{\eqref{eq:loc-maps}}{\cong}[\qfiveEP\otimes N^T S^1,\Sigma^k HM]^G
=[\qfiveEP\otimes N^T S^1,\qfiveEP\otimes\Sigma^k HM]^G
\overset{\eqref{eq:phi-maps}}{\cong}[\Phi^G N^T S^1,\Phi^G\Sigma^k HM]_{\qfiveSp},
\]
where in the middle equality we used $\Sigma^k HM\simeq\qfiveEP\otimes\Sigma^k HM$. By Proposition~\ref{prop:phiNT}, $\Phi^G N^T S^1\simeq S^{|T/G|}$, while
$\Phi^G\Sigma^k HM=\Sigma^k HA$ by (G2). Hence
\[
[N^T S^1,\Sigma^k HM]^G\cong[S^{|T/G|},\Sigma^k HA]_{\qfiveSp}
=\pi_{|T/G|-k}(HA)=0 ,
\]
since $|T/G|-k\ge 1>0$ and $HA$ is concentrated in degree $0$.
\end{proof}

\begin{qfivecorollary}\label{cor:exact-slice}
For geometric $M$ and $k\ge0$, $\Sigma^k HM$ is an $\qfiveaO$-$(k r_G)$-slice, where
$r_G=[G:\chi^\qfiveaO(G)]$.
\end{qfivecorollary}

\begin{proof}
\emph{$(kr_G)$-connectivity.} By (G1), $\Sigma^k HM\simeq\qfiveEP\otimes\Sigma^k HM$,
which lies in $\tau^\qfiveaO_{\ge k r_G}$ by Lemma \ref{lem:geom}(2); thus
$\Sigma^k HM$ is $\qfiveaO$-slice $(k r_G)$-connective.

\emph{Upper truncation: what must be shown.} The $\qfiveaO$-slice $(k r_G)$-truncation
$P^{k r_G}_\qfiveaO$ is the nullification killing $\tau^\qfiveaO_{\ge k r_G+1}$. By its
universal property, the natural map $\Sigma^kHM\to P^{kr_G}_\qfiveaO\Sigma^kHM$ is an
equivalence if and only if $\Sigma^kHM$ is \emph{$\tau^\qfiveaO_{\ge kr_G+1}$-null},
i.e.\ $[X,\Sigma^kHM]^G=0$ for \emph{every} $X\in\tau^\qfiveaO_{\ge kr_G+1}$. Write
$W:=\Sigma^kHM$ and recall
\[
\tau^\qfiveaO_{\ge kr_G+1}=\langle\mathcal G\rangle_{\ge},\qquad
\mathcal G:=\bigl\{\,G_+\otimes_H N^T S^1\ \bigm|\ H\subseteq G,\ T\in\qfiveaO(H),\ |T|>k r_G\,\bigr\}.
\]

\emph{The closure class.} Define
\[
\mathcal N_W:=\bigl\{\,X\in\qfiveSp^G\ \bigm|\
[\Sigma^\ell(G/K_+\otimes X),W]^G=0\ \text{ for all }K\subseteq G\text{ and all }\ell\ge0\,\bigr\}.
\]
(Taking $K=G$ and $\ell=0$ shows $X\in\mathcal N_W\Rightarrow[X,W]^G=0$, so
$\tau^\qfiveaO_{\ge kr_G+1}\subseteq\mathcal N_W$ will give exactly the desired
nullity.) We claim $\mathcal N_W$ is closed under the generation operations
(o1)--(o3) and contains $\mathcal G$; then
$\tau^\qfiveaO_{\ge kr_G+1}=\langle\mathcal G\rangle_{\ge}\subseteq\mathcal N_W$, as
required.

\emph{Closure of $\mathcal N_W$.} Fix $K\subseteq G$ and $\ell\ge0$, and note the
functor $[\Sigma^\ell(G/K_+\otimes(-)),W]^G$ is the $\ell$th homotopy group of the
mapping spectrum $\operatorname{Map}^G(G/K_+\otimes(-),W)$, hence sends cofiber
sequences in the source to long exact sequences.
\begin{itemize}\setlength\itemsep{2pt}
\item \emph{(o1) Positive suspensions, colimits, cofibers, extensions.} For
$X\in\mathcal N_W$ and $j\ge0$, $[\Sigma^\ell(G/K_+\otimes\Sigma^jX),W]^G
=[\Sigma^{\ell+j}(G/K_+\otimes X),W]^G=0$ since $\ell+j\ge0$; thus $\mathcal N_W$
is closed under positive suspension. For a cofiber sequence $X\to Y\to Z$ with
$X,Y\in\mathcal N_W$, the long exact sequence
\[
[\Sigma^{\ell+1}(G/K_+\otimes X),W]^G\to[\Sigma^\ell(G/K_+\otimes Z),W]^G
\to[\Sigma^\ell(G/K_+\otimes Y),W]^G
\]
has vanishing outer terms for $\ell\ge0$ (as $\ell+1,\ell\ge0$), forcing the
middle term to vanish; so $Z\in\mathcal N_W$. Wedges and filtered colimits send
$\operatorname{Map}^G(G/K_+\otimes(-),W)$ to products and homotopy limits; for a
homotopy limit of spectra each having vanishing $\pi_{\ge0}$ the $\lim^1$-sequence
$0\to{\lim}^1\pi_{\ell+1}\to\pi_\ell\to\lim\pi_\ell\to0$ gives vanishing $\pi_\ell$
for $\ell\ge0$ (both $\ell,\ell+1\ge0$). Hence $\mathcal N_W$ is closed under all
homotopy colimits, and therefore under extensions.
\item \emph{(o2) Retracts.} A retract of an object with vanishing
$[\Sigma^\ell(G/K_+\otimes-),W]^G$ has the same vanishing.
\item \emph{(o3) Orbit tensors.} Let $X\in\mathcal N_W$ and $K'\subseteq G$.
Using $G/K_+\otimes G/K'_+\simeq\bigsqcup_{[g]\in K\backslash G/K'}
G/(K\cap gK'g^{-1})_+$, we get for each $K,\ell$
\[
[\Sigma^\ell(G/K_+\otimes(G/K'_+\otimes X)),W]^G
\cong\!\!\prod_{[g]\in K\backslash G/K'}\!\![\Sigma^\ell(G/(K\cap gK'g^{-1})_+\otimes X),W]^G=0,
\]
each factor vanishing because $X\in\mathcal N_W$ (with the subgroup
$K\cap gK'g^{-1}$ in the r\^ole of $K$). Hence $G/K'_+\otimes X\in\mathcal N_W$.
\end{itemize}

\emph{$\mathcal G\subseteq\mathcal N_W$.} Let $C=G_+\otimes_H N^T S^1\in\mathcal G$
($T\in\qfiveaO(H)$, $|T|>kr_G$); we must show $[\Sigma^\ell(G/K_+\otimes C),W]^G=0$
for all $K\subseteq G$, $\ell\ge0$. By the projection formula and the
double-coset norm--restriction formula \eqref{eq:resnorm},
\[
G/K_+\otimes C\simeq G_+\otimes_H\bigl(i_H^\ast(G/K_+)\otimes N^TS^1\bigr)
\simeq\bigsqcup_{j} G_+\otimes_{L_j} N^{\,i_{L_j}^\ast T}S^1,
\]
where $i_H^\ast(G/K)=\bigsqcup_j H/L_j$ and we used $(H/L_j)_+\otimes N^TS^1
\simeq H_+\otimes_{L_j}N^{i_{L_j}^\ast T}S^1$. Each $i_{L_j}^\ast T$ is admissible
for $\qfiveaO|_{L_j}$ (admissibility is closed under restriction) and has the same
underlying cardinality $|i_{L_j}^\ast T|=|T|>kr_G$; hence each summand
$G_+\otimes_{L_j}N^{i_{L_j}^\ast T}S^1$ again belongs to $\mathcal G$. Thus it
suffices to prove $[\Sigma^\ell C',W]^G=0$ for $\ell\ge0$ and every
\emph{generator} $C'=G_+\otimes_{H'}N^{T'}S^1\in\mathcal G$ (with $|T'|>kr_G$),
since $[\Sigma^\ell(G/K_+\otimes C),W]^G=\prod_j[\Sigma^\ell C'_j,W]^G$.

Fix such a generator $C'=G_+\otimes_{H'}N^{T'}S^1$ and $\ell\ge0$. By the
suspension identity $\Sigma^\ell(G_+\otimes_{H'}N^{T'}S^1)\simeq
G_+\otimes_{H'}N^{T'\sqcup\ell\ast}S^1$ (with $\ell\ast$ a disjoint union of $\ell$ trivial $H'$-orbits $H'/H'$, so
$T'\sqcup\ell\ast\in\qfiveaO(H')$ of cardinality $|T'|+\ell>kr_G$) and the
induction/restriction adjunction,
\[
[\Sigma^\ell C',W]^G\cong[N^{\,T'\sqcup\ell\ast}S^1,\,i_{H'}^\ast W]^{H'},
\qquad W=\Sigma^kHM.
\]
\emph{Case $H'\subsetneq G$ proper.} We claim $i_{H'}^\ast HM\simeq 0$.
Indeed, restriction $i_{H'}^\ast$ is exact and preserves the homotopy
$t$-structure, so $i_{H'}^\ast HM$ is again an Eilenberg--MacLane $H'$-spectrum,
i.e.\ its homotopy is concentrated in degree $0$ with
$\underline\pi_0(i_{H'}^\ast HM)=i_{H'}^\ast M$, the restricted Mackey functor.
For the geometric Mackey functor $M$ one has, for every $J\subseteq H'$,
\[
(i_{H'}^\ast M)(H'/J)\;=\;M(G/J)\;=\;0,
\]
because $J\subseteq H'\subsetneq G$ is a proper subgroup of $G$ and $M(G/J)=0$ by
(M1) (this is exactly the computation carried out in the proof of
Proposition~\ref{prop:geomEM}(G1)). Thus $i_{H'}^\ast HM$ is an
Eilenberg--MacLane spectrum on the zero Mackey functor, hence contractible:
$i_{H'}^\ast HM\simeq 0$. Therefore
$i_{H'}^\ast W=\Sigma^k i_{H'}^\ast HM\simeq 0$ and the group
$[N^{\,T'\sqcup\ell\ast}S^1,i_{H'}^\ast W]^{H'}=0$ vanishes \emph{outright},
without reference to any $H'$-level connectivity or geometric threshold.
\emph{This is the crucial point flagged in the issue: at proper $H'$ we do not
invoke Lemma~\ref{lem:geom}(3) over $H'$ (which would only control degrees
$|T|>kr_{H'}$ and is moreover unavailable since $i_{H'}^\ast HM$ need not be
geometric over $H'$); we use the strictly stronger fact that the target itself
restricts to $0$. Consequently the potential degree mismatch $r_{H'}<r_G$ never
enters.} \emph{Case $H'=G$.} Then $T'\sqcup\ell\ast\in\qfiveaO(G)$ has cardinality
$>kr_G$, and Lemma \ref{lem:geom}(3), applied to this admissible $G$-set, gives
\[
[N^{\,T'\sqcup\ell\ast}S^1,\Sigma^kHM]^G\cong
\pi_{\,|(T'\sqcup\ell\ast)/G|-k}(HA)=0,
\]
since $|(T'\sqcup\ell\ast)/G|=|T'/G|+\ell\ge(k+1)+\ell>k$ (using $|T'|>kr_G
\Rightarrow|T'/G|>k$ by Lemma \ref{lem:index}(3)) and $HA$ is concentrated in
degree $0$. In both cases $[\Sigma^\ell C',W]^G=0$ for all $\ell\ge0$, so
$\mathcal G\subseteq\mathcal N_W$.

\emph{Conclusion.} By the two claims,
$\tau^\qfiveaO_{\ge kr_G+1}=\langle\mathcal G\rangle_{\ge}\subseteq\mathcal N_W$, and
in particular (taking $K=G$, $\ell=0$) $[X,\Sigma^kHM]^G=0$ for every
$X\in\tau^\qfiveaO_{\ge kr_G+1}$. Hence $\Sigma^kHM$ is $\tau^\qfiveaO_{\ge kr_G+1}$-null
and $\Sigma^kHM\to P^{kr_G}_\qfiveaO\Sigma^kHM$ is an equivalence. Together with
$\Sigma^kHM\in\tau^\qfiveaO_{\ge kr_G}$ from the first paragraph, $\Sigma^k HM$ is
exactly a $(k r_G)$-slice.
\end{proof}

\begin{proof}[Proof of $P(G,\qfiveaO)$ (converse), completing Theorem \ref{thm:main}]
Let $E$ be connective with $[H:\chi^\qfiveaO(H)]\cdot\qfivegconn(E)(H)\ge n$ for all
$H\subseteq G$. Consider the isotropy-separation cofiber sequence
\[
E\mathcal P_+\otimes E\;\longrightarrow\; E\;\longrightarrow\;\qfiveEP\otimes E .
\]
As $\tau^\qfiveaO_{\ge n}$ is closed under cofibers/extensions, it suffices to show
both ends lie in $\tau^\qfiveaO_{\ge n}$.

\emph{The induced part.} For each proper $H$, Lemma \ref{lem:restr}(1)-(2) shows
$i_H^\ast E$ satisfies, for all $K\subseteq H$,
\[
[K:\chi^{\qfiveaO|_H}(K)]\cdot\qfivegconn(i_H^\ast E)(K)
=[K:\chi^\qfiveaO(K)]\cdot\qfivegconn(E)(K)\ge n ,
\]
and $i_H^\ast E$ is connective. By the inductive hypothesis $P(H,\qfiveaO|_H)$ (valid
since $H$ is proper), $i_H^\ast E\in\tau^{\qfiveaO|_H}_{\ge n}$. By Lemma
\ref{lem:induced}, $E\mathcal P_+\otimes E\in\tau^\qfiveaO_{\ge n}$.

\emph{The geometric part.} Set $Y:=\qfiveEP\otimes E$, an $\qfiveEP$-local spectrum with
$\Phi^G Y\simeq\Phi^G E$ (by (L2)), of connectivity $c:=\qfivegconn(E)(G)$. If
$c=+\infty$ (i.e.\ $\Phi^G E\simeq 0$) then $Y\simeq 0\in\tau^\qfiveaO_{\ge n}$ and we are
done; assume $c<\infty$, so $c\ge 0$ (because $E$ is connective, $\Phi^G E$ is
connective, $c\ge0$). The hypothesis at $H=G$ gives $r_G\,c\ge n$.

The spectrum $\Phi^G E$ is $c$-connective, hence admits a CW/Whitehead filtration:
it is a sequential colimit $\Phi^G E\simeq\operatorname*{colim}_q F_q$ where
$F_c=\Sigma^c H\pi_c(\Phi^G E)$ and each $F_{q}\to F_{q+1}$ ($q\ge c$) fits in a
cofiber sequence
\[
F_q\longrightarrow F_{q+1}\longrightarrow \Sigma^{q+1}H\pi_{q+1}(\Phi^G E),
\]
the $(q+1)$-st Postnikov stage \cite[\S2]{HillPrimer}. Apply the inverse
equivalence $\Phi^{-1}:\qfiveSp\xrightarrow{\sim}L\,\qfiveSp^G$ of (L2). Since $\Phi^{-1}$
preserves colimits and cofiber sequences and $Y=\Phi^{-1}(\Phi^G E)$, we get
$Y\simeq\operatorname*{colim}_q\Phi^{-1}F_q$ with cofiber sequences
\[
\Phi^{-1}F_q\longrightarrow\Phi^{-1}F_{q+1}\longrightarrow
\Phi^{-1}\bigl(\Sigma^{q+1}H\pi_{q+1}(\Phi^G E)\bigr).
\]
For each abelian group $A_j:=\pi_j(\Phi^G E)$ ($j\ge c$), let
$M_j:=\underline{A_j}^{\mathrm{geo}}$ be the geometric Mackey functor with
geometric value $A_j$ (Definition \ref{def:geom}); by Proposition
\ref{prop:geomEM} its Eilenberg--MacLane spectrum $HM_j$ is connective,
$\qfiveEP$-local, and satisfies $\Phi^G HM_j\simeq HA_j$. Then
$\Phi^{-1}(\Sigma^j HA_j)\simeq\qfiveEP\otimes\Sigma^j HM_j$ (both are $\qfiveEP$-local with
geometric fixed points $\Sigma^j HA_j$, so equal by (L2)). By Lemma
\ref{lem:geom}(2) (whose connectivity hypothesis is now verified for $M_j$),
\[
\qfiveEP\otimes\Sigma^j HM_j\in\tau^\qfiveaO_{\ge j r_G}\subseteq\tau^\qfiveaO_{\ge c r_G}
\subseteq\tau^\qfiveaO_{\ge n}\qquad(j\ge c,\ \ j r_G\ge c r_G\ge n).
\]
Thus $\Phi^{-1}F_c=\qfiveEP\otimes\Sigma^c HM_c\in\tau^\qfiveaO_{\ge n}$, and each layer of
the filtration of $Y$ lies in $\tau^\qfiveaO_{\ge n}$; since $\tau^\qfiveaO_{\ge n}$ is closed
under extensions and sequential colimits, $Y=\qfiveEP\otimes E\in\tau^\qfiveaO_{\ge n}$.

Both ends being in $\tau^\qfiveaO_{\ge n}$, we conclude $E\in\tau^\qfiveaO_{\ge n}$, proving
$P(G,\qfiveaO)$ and, by induction, the converse for all groups.
\end{proof}

This is precisely the $\qfiveaO$-relative form of \cite[Thm.~2.5]{HY18}; the only
new ingredient is that the multiplier $|H|$ is replaced by $[H:\chi^\qfiveaO(H)]$,
reflecting that $\qfiveaO$ admits norms only up to index $[H:\chi^\qfiveaO(H)]$ per orbit,
encoded by the degree-shift Lemma \ref{lem:shift}.

\section*{Examples and consistency checks}

\paragraph{Complete $\qfiveaO$.} Here every $G/H$ is admissible, $\chi^\qfiveaO(H)=\{e\}$,
$r_H=|H|$, and \eqref{eq:crit} reads $|H|\cdot\qfivegconn(E)(H)\ge n$ for all
$H$ --- the Hill--Hopkins--Ravenel slice connectivity in the Hill--Yarnall
formulation \cite{HY18, HHR16}. The cell $C_G=N^{G/e}S^1=S^{\rho_G}$ is the
regular-representation sphere; Lemma \ref{lem:shift} is then the classical fact
that smashing with $S^{\rho_G}$ raises slice degree by $|G|$.

\paragraph{Trivial $\qfiveaO$.} No nontrivial transfers: $\chi^\qfiveaO(H)=H$ for all $H$,
so $r_H=1$ and \eqref{eq:crit} becomes $\qfivegconn(E)(H)\ge n$ for all
$H$. The only admissible sets are trivial orbits ($C_G=N^{G/G}S^1=S^1$,
$r_G=1$), and the filtration is the geometric Postnikov tower; Lemma
\ref{lem:shift} reduces to ordinary suspension.

\paragraph{$G=\qfiveZ/p$.} The two transfer systems are trivial and complete. For the
complete one, $\chi^\qfiveaO(\qfiveZ/p)=\{e\}$ gives the bound $p\cdot\qfivegconn(E)(\qfiveZ/p)\ge n$
at the top together with $\qfivegconn(E)(e)\ge n$; for the trivial one both indices
are $1$ and the conditions are $\qfivegconn(E)(H)\ge n$. These agree with the
$C_p$-slice computations of \cite{HY18}.

\section*{Scope}

The model-categorical and $\infty$-categorical foundations for the
norm functors, the projection formula, the geometric-fixed-point and double-coset
formulas \eqref{eq:phinorm}--\eqref{eq:resnorm}, and equivariant localization
(L1)--(L2) are taken from \cite{HHR16, HillPrimer, BBR, Rubin}. The argument is the
Hill--Yarnall template \cite{HY18} carried out with the
indexing-system-restricted norms in place of the full regular-representation
cells; the new technical inputs are the degree-shift Lemma \ref{lem:shift} and the
geometric-slice Lemma \ref{lem:geom}, which together replace the previously
unproved equivalence and make the characteristic function $\chi^\qfiveaO$ govern the
slice degree.

\begin{thebibliography}{9}
\bibitem{HHR16} M.~A.~Hill, M.~J.~Hopkins, D.~C.~Ravenel.
  \emph{On the nonexistence of elements of Kervaire invariant one}.
  Ann.\ of Math.\ (2) \textbf{184} (2016), 1--262.
\bibitem{HY18} M.~A.~Hill, C.~Yarnall.
  \emph{A new formulation of the equivariant slice filtration with applications
  to $C_p$-slices}. Proc.\ Amer.\ Math.\ Soc.\ \textbf{146} (2018), 3605--3614.
\bibitem{HillPrimer} M.~A.~Hill.
  \emph{The equivariant slice filtration: a primer}.
  Homology Homotopy Appl.\ \textbf{14} (2012), 143--166.
\bibitem{BBR} S.~Balchin, D.~Barnes, C.~Roitzheim.
  \emph{$N_\infty$-operads and associahedra}.
  Pacific J.\ Math.\ \textbf{315} (2021), 285--304.
\bibitem{Rubin} J.~Rubin.
  \emph{Detecting Steiner and linear isometries operads}.
  Glasg.\ Math.\ J.\ \textbf{63} (2021), 307--342.
\bibitem{Wilson} D.~Wilson.
  \emph{On categories of slices}. Preprint, arXiv:1711.03472, 2017.
\end{thebibliography}

\endgroup